\color{unidarkgreen}{\section[(обязательное) Уставки функций РЗА]{(обязательное)\\Уставки функций РЗА}\label{app:set}}
\color{black}


%%%%%%%%%%%%%%%%%%%%%%%%%%%%%%%%%%%%%%%%%%%%%%%%%%%%%%%%%% ДЗ  %%%%%%%%%%%%%%%%%%%%%%%%%%%%%%%%%%%%%%%%%%%%%%%%%%%%%%%%%%%%%%%%%%
\par
Параметры для настройки ДЗ приведены в таблице \ref{dz:tbl0}.

\footnotesize
\begin{longtable}{|>{\centering\arraybackslash}m{5.3cm}|>{\centering\arraybackslash}m{3.3cm}|>{\centering\arraybackslash}m{4.2cm}|>{\centering\arraybackslash}m{1.8cm}|>{\centering\arraybackslash}m{1cm}|}
\caption{Общие параметры для настройки ДЗ\hfill\vspace{-0.5\baselineskip}}\label{dz:tbl0}\\ 
\hline
\rowcolor{gray!20}
Параметр (Параметр на ИЧМ) & Условное обозначение на схеме & Значение/ Диапазон & Единица измерения & Шаг \\ 
\hline
\endfirsthead
\caption*{\hspace{3pt}\emph{Продолжение таблицы \ref{dz:tbl0}\hfill\vspace{-0.5\baselineskip}}} \\ % сделано по ГОСТ 2.105 п.6.8.7
\hline
\rowcolor{gray!20}
Параметр (Параметр на ИЧМ) & Условное обозначение на схеме & Значение/ Диапазон & Единица измерения & Шаг \\ 
%\hline
\endhead
%\hline
\endfoot
%\hline
\endlastfoot
%==+t1*LLN0|LVDIS
\multicolumn{5}{|c|}{ ДЗ-фф 1 ст } \\ \hline 
\centering Ввод функции в работу (Ввод\_функции)  & \centering SGF1 & \centering Не предусмотрено \\ Предусмотрено & \centering -- & \centering \arraybackslash -- \\
\hline
\centering Угол наклона верхней стороны характеристики срабатывания для компенсации нагрузки (Фкн)  & \centering -- & \centering -45,0...0,0 & \centering град. & \centering \arraybackslash 0,1 \\
\hline
\centering Активное сопротивление срабатывания (Rср)  & \centering -- & \centering 0,01...500,00 & \centering Ом & \centering \arraybackslash 0,01 \\
\hline
\centering Реактивное сопротивление срабатывания (Xср)  & \centering -- & \centering 0,01...500,00 & \centering Ом & \centering \arraybackslash 0,01 \\
\hline
\centering Угол наклона характеристики срабатывания (Фхс)  & \centering -- & \centering 30,0...89,9 & \centering град. & \centering \arraybackslash 0,1 \\
\hline
\centering Режим направленности (Направленность)  & \centering SGF3 & \centering Ненаправленная \\ Прямонаправленная \\ Обратнонаправленная & \centering -- & \centering \arraybackslash -- \\
\hline
\centering Подхват от ненаправленного пуска или отдельной ступени (Подхват)  & \centering SGF9 & \centering Не предусмотрено \\ Предусмотрено & \centering -- & \centering \arraybackslash -- \\
\hline
\centering Режим контроля от БК (Реж\_БК)  & \centering SGF7 & \centering С контролем от БК-б \\ С контролем от БК-м & \centering -- & \centering \arraybackslash -- \\
\hline
\centering Режим контроля от БНН (Реж\_БНН)  & \centering SGF8 & \centering Без контроля \\ С контролем & \centering -- & \centering \arraybackslash -- \\
\hline
\centering Режим контроля от ОПО (Реж\_ОПО)  & \centering SGF6 & \centering По току \\ По току и напряжению \\ По току или по току и напряжению & \centering -- & \centering \arraybackslash -- \\
\hline
\centering Тип ПО (Контр\_ПО)  & \centering SGF5 & \centering Без контроля \\ С контролем от БК \\ С контролем от ОПО & \centering -- & \centering \arraybackslash -- \\
\hline
\centering Вывод направленности при АУ (Ненапр\_АУ)  & \centering SGF4 & \centering Не предусмотрено \\ Предусмотрено & \centering -- & \centering \arraybackslash -- \\
\hline
\centering Выдержка времени срабатывания (Tср)  & \centering T1 & \centering 0...10000 & \centering мс & \centering \arraybackslash 10 \\
\hline
\multicolumn{5}{|c|}{ ДЗ-фф 2 ст } \\ \hline 
\centering Ввод функции в работу (Ввод\_функции)  & \centering SGF1 & \centering Не предусмотрено \\ Предусмотрено & \centering -- & \centering \arraybackslash -- \\
\hline
\centering Угол наклона верхней стороны характеристики срабатывания для компенсации нагрузки (Фкн)  & \centering -- & \centering -45,0...0,0 & \centering град. & \centering \arraybackslash 0,1 \\
\hline
\centering Активное сопротивление срабатывания (Rср)  & \centering -- & \centering 0,01...500,00 & \centering Ом & \centering \arraybackslash 0,01 \\
\hline
\centering Реактивное сопротивление срабатывания (Xср)  & \centering -- & \centering 0,01...500,00 & \centering Ом & \centering \arraybackslash 0,01 \\
\hline
\centering Угол наклона характеристики срабатывания (Фхс)  & \centering -- & \centering 30,0...89,9 & \centering град. & \centering \arraybackslash 0,1 \\
\hline
\centering Режим направленности (Направленность)  & \centering SGF3 & \centering Ненаправленная \\ Прямонаправленная \\ Обратнонаправленная & \centering -- & \centering \arraybackslash -- \\
\hline
\centering Режим контроля от БК (Реж\_БК)  & \centering SGF7 & \centering С контролем от БК-б \\ С контролем от БК-м & \centering -- & \centering \arraybackslash -- \\
\hline
\centering Режим контроля от БНН (Реж\_БНН)  & \centering SGF8 & \centering Без контроля \\ С контролем & \centering -- & \centering \arraybackslash -- \\
\hline
\centering Режим контроля от ОПО (Реж\_ОПО)  & \centering SGF6 & \centering По току \\ По току и напряжению \\ По току или по току и напряжению & \centering -- & \centering \arraybackslash -- \\
\hline
\centering Тип ПО (Контр\_ПО)  & \centering SGF5 & \centering Без контроля \\ С контролем от БК \\ С контролем от ОПО & \centering -- & \centering \arraybackslash -- \\
\hline
\centering Вывод направленности при АУ (Ненапр\_АУ)  & \centering SGF4 & \centering Не предусмотрено \\ Предусмотрено & \centering -- & \centering \arraybackslash -- \\
\hline
\centering Выдержка времени срабатывания (Tср)  & \centering T1 & \centering 0...10000 & \centering мс & \centering \arraybackslash 10 \\
\hline
\multicolumn{5}{|c|}{ ДЗ-фф 3 ст } \\ \hline 
\centering Ввод функции в работу (Ввод\_функции)  & \centering SGF1 & \centering Не предусмотрено \\ Предусмотрено & \centering -- & \centering \arraybackslash -- \\
\hline
\centering Угол наклона верхней стороны характеристики срабатывания для компенсации нагрузки (Фкн)  & \centering -- & \centering -45,0...0,0 & \centering град. & \centering \arraybackslash 0,1 \\
\hline
\centering Активное сопротивление срабатывания (Rср)  & \centering -- & \centering 0,01...500,00 & \centering Ом & \centering \arraybackslash 0,01 \\
\hline
\centering Реактивное сопротивление срабатывания (Xср)  & \centering -- & \centering 0,01...500,00 & \centering Ом & \centering \arraybackslash 0,01 \\
\hline
\centering Угол наклона характеристики срабатывания (Фхс)  & \centering -- & \centering 30,0...89,9 & \centering град. & \centering \arraybackslash 0,1 \\
\hline
\centering Режим направленности (Направленность)  & \centering SGF3 & \centering Ненаправленная \\ Прямонаправленная \\ Обратнонаправленная & \centering -- & \centering \arraybackslash -- \\
\hline
\centering Режим контроля от БК (Реж\_БК)  & \centering SGF7 & \centering С контролем от БК-б \\ С контролем от БК-м & \centering -- & \centering \arraybackslash -- \\
\hline
\centering Режим контроля от БНН (Реж\_БНН)  & \centering SGF8 & \centering Без контроля \\ С контролем & \centering -- & \centering \arraybackslash -- \\
\hline
\centering Режим контроля от ОПО (Реж\_ОПО)  & \centering SGF6 & \centering По току \\ По току и напряжению \\ По току или по току и напряжению & \centering -- & \centering \arraybackslash -- \\
\hline
\centering Тип ПО (Контр\_ПО)  & \centering SGF5 & \centering Без контроля \\ С контролем от БК \\ С контролем от ОПО & \centering -- & \centering \arraybackslash -- \\
\hline
\centering Вывод направленности при АУ (Ненапр\_АУ)  & \centering SGF4 & \centering Не предусмотрено \\ Предусмотрено & \centering -- & \centering \arraybackslash -- \\
\hline
\centering Выдержка времени срабатывания (Tср)  & \centering T1 & \centering 0...10000 & \centering мс & \centering \arraybackslash 10 \\
\hline
\multicolumn{5}{|c|}{ ДЗ-фз 1 ст } \\ \hline 
\centering Ввод функции в работу (Ввод\_функции)  & \centering SGF1 & \centering Не предусмотрено \\ Предусмотрено & \centering -- & \centering \arraybackslash -- \\
\hline
\centering Угол наклона верхней стороны характеристики срабатывания для компенсации нагрузки (Фкн)  & \centering -- & \centering -45,0...0,0 & \centering град. & \centering \arraybackslash 0,1 \\
\hline
\centering Активное сопротивление срабатывания (Rср)  & \centering -- & \centering 0,01...500,00 & \centering Ом & \centering \arraybackslash 0,01 \\
\hline
\centering Реактивное сопротивление срабатывания (Xср)  & \centering -- & \centering 0,01...500,00 & \centering Ом & \centering \arraybackslash 0,01 \\
\hline
\centering Угол наклона характеристики срабатывания (Фхс)  & \centering -- & \centering 30,0...89,9 & \centering град. & \centering \arraybackslash 0,1 \\
\hline
\centering Режим направленности (Направленность)  & \centering SGF3 & \centering Ненаправленная \\ Прямонаправленная \\ Обратнонаправленная & \centering -- & \centering \arraybackslash -- \\
\hline
\centering Подхват от ненаправленного пуска или отдельной ступени (Подхват)  & \centering SGF9 & \centering Не предусмотрено \\ Предусмотрено & \centering -- & \centering \arraybackslash -- \\
\hline
\centering Режим контроля от БК (Реж\_БК)  & \centering SGF7 & \centering С контролем от БК-б \\ С контролем от БК-м & \centering -- & \centering \arraybackslash -- \\
\hline
\centering Режим контроля от БНН (Реж\_БНН)  & \centering SGF8 & \centering Без контроля \\ С контролем & \centering -- & \centering \arraybackslash -- \\
\hline
\centering Режим контроля от ОПО (Реж\_ОПО)  & \centering SGF6 & \centering По току \\ По току и напряжению \\ По току или по току и напряжению & \centering -- & \centering \arraybackslash -- \\
\hline
\centering Тип ПО (Контр\_ПО)  & \centering SGF5 & \centering Без контроля \\ С контролем от БК \\ С контролем от ОПО & \centering -- & \centering \arraybackslash -- \\
\hline
\centering Вывод направленности при АУ (Ненапр\_АУ)  & \centering SGF4 & \centering Не предусмотрено \\ Предусмотрено & \centering -- & \centering \arraybackslash -- \\
\hline
\centering Выдержка времени срабатывания (Tср)  & \centering T1 & \centering 0...10000 & \centering мс & \centering \arraybackslash 10 \\
\hline
\multicolumn{5}{|c|}{ ДЗ-фз 2 ст } \\ \hline 
\centering Ввод функции в работу (Ввод\_функции)  & \centering SGF1 & \centering Не предусмотрено \\ Предусмотрено & \centering -- & \centering \arraybackslash -- \\
\hline
\centering Угол наклона верхней стороны характеристики срабатывания для компенсации нагрузки (Фкн)  & \centering -- & \centering -45,0...0,0 & \centering град. & \centering \arraybackslash 0,1 \\
\hline
\centering Активное сопротивление срабатывания (Rср)  & \centering -- & \centering 0,01...500,00 & \centering Ом & \centering \arraybackslash 0,01 \\
\hline
\centering Реактивное сопротивление срабатывания (Xср)  & \centering -- & \centering 0,01...500,00 & \centering Ом & \centering \arraybackslash 0,01 \\
\hline
\centering Угол наклона характеристики срабатывания (Фхс)  & \centering -- & \centering 30,0...89,9 & \centering град. & \centering \arraybackslash 0,1 \\
\hline
\centering Режим направленности (Направленность)  & \centering SGF3 & \centering Ненаправленная \\ Прямонаправленная \\ Обратнонаправленная & \centering -- & \centering \arraybackslash -- \\
\hline
\centering Режим контроля от БК (Реж\_БК)  & \centering SGF7 & \centering С контролем от БК-б \\ С контролем от БК-м & \centering -- & \centering \arraybackslash -- \\
\hline
\centering Режим контроля от БНН (Реж\_БНН)  & \centering SGF8 & \centering Без контроля \\ С контролем & \centering -- & \centering \arraybackslash -- \\
\hline
\centering Режим контроля от ОПО (Реж\_ОПО)  & \centering SGF6 & \centering По току \\ По току и напряжению \\ По току или по току и напряжению & \centering -- & \centering \arraybackslash -- \\
\hline
\centering Тип ПО (Контр\_ПО)  & \centering SGF5 & \centering Без контроля \\ С контролем от БК \\ С контролем от ОПО & \centering -- & \centering \arraybackslash -- \\
\hline
\centering Вывод направленности при АУ (Ненапр\_АУ)  & \centering SGF4 & \centering Не предусмотрено \\ Предусмотрено & \centering -- & \centering \arraybackslash -- \\
\hline
\centering Выдержка времени срабатывания (Tср)  & \centering T1 & \centering 0...10000 & \centering мс & \centering \arraybackslash 10 \\
\hline
\multicolumn{5}{|c|}{ АУ ДЗ } \\ \hline 
\centering Выбор ускоряемой ступени (Выбор\_ст)  & \centering SGF1 & \centering 2 ступень \\ 3 ступень & \centering -- & \centering \arraybackslash -- \\
\hline
\multicolumn{5}{|c|}{ ОПО } \\ \hline 
\centering Ток срабатывания ПО ДЗ по току (Iср)  & \centering I\verb|>>| & \centering 0,10...5,00 & \centering о.е. & \centering \arraybackslash 0,01 \\
\hline
\centering Ток срабатывания ПО ДЗ по току и напряжению (IсрU)  & \centering I> & \centering 0,10...5,00 & \centering о.е. & \centering \arraybackslash 0,01 \\
\hline
\centering Напряжение срабатывания ПО по току и напряжению (Uср)  & \centering U< & \centering 2,0...100,0 & \centering В & \centering \arraybackslash 0,1 \\
\hline
\centering Режим контроля от БНН (Реж\_БНН)  & \centering SGF1 & \centering Не предусмотрено \\ Предусмотрено & \centering -- & \centering \arraybackslash -- \\
\hline
\multicolumn{5}{|c|}{ ОВП } \\ \hline 
\centering Коэффициент отношения тока нулевой последовательности к току прямой последовательности (Kдел.)  & \centering 3I0/I1> & \centering 0,10...2,00 & \centering о.е. & \centering \arraybackslash 0,01 \\
\hline
\centering Уставка по напряжению 3U0 (3U0ср)  & \centering 3U0> & \centering 2,0...100,0 & \centering В & \centering \arraybackslash 0,1 \\
\hline
\centering Напряжение срабатывания (Uср)  & \centering U< & \centering 2,0...100,0 & \centering В & \centering \arraybackslash 0,1 \\
\hline
\centering Режим контроля от БНН (Реж\_БНН)  & \centering SGF1 & \centering Не предусмотрено \\ Предусмотрено & \centering -- & \centering \arraybackslash -- \\
\hline
\multicolumn{5}{|c|}{ ОН } \\ \hline 
\centering Угол направленности во II квадранте (Фнапр\_II)  & \centering φ2 & \centering 90,0...130,0 & \centering град. & \centering \arraybackslash 0,1 \\
\hline
\centering Угол направленности в IV квадранте (Фнапр\_IV)  & \centering φ1 & \centering -40,0...0,0 & \centering град. & \centering \arraybackslash 0,1 \\
\hline
\multicolumn{5}{|c|}{ ОУ ДЗ } \\ \hline 
\centering Выбор оперативно ускоряемой ступени (Выбор\_ст\_ОУ)  & \centering SGF1 & \centering 2 ступень \\ 3 ступень & \centering -- & \centering \arraybackslash -- \\
\hline
\centering Выдержка времени срабатывания оперативного ускорения (T\_ОУ)  & \centering T1 & \centering 0...3000 & \centering мс & \centering \arraybackslash 10 \\
\hline
\multicolumn{5}{|c|}{ ДЗ } \\ \hline 
\centering Учет нагрузочного режима в характеристике срабатывания (Вырез\_нагрузки)  & \centering SGF2 & \centering Не предусмотрено \\ Предусмотрено & \centering -- & \centering \arraybackslash -- \\
\hline
\centering Активное сопротивление для отстройки от режима максимальной нагрузки (Rнг)  & \centering -- & \centering 0,01...500,00 & \centering Ом & \centering \arraybackslash 0,01 \\
\hline
\centering Угол наклона для отстройки от режима максимальной нагрузки (Фнг)  & \centering -- & \centering 5,0...60,0 & \centering град. & \centering \arraybackslash 0,1 \\
\hline
\centering Режим работы ДЗ по контуру "фаза В - земля" (Ввод\_B0)  & \centering SGF10 & \centering Не предусмотрено \\ Предусмотрено & \centering -- & \centering \arraybackslash -- \\
\hline
\centering Удельное активное сопротивление прямой последовательности ()  & \centering R1уд & \centering 0,001...100,000 & \centering Ом/км & \centering \arraybackslash 0,001 \\
\hline
\centering Удельное реактивное сопротивление прямой последовательности ()  & \centering X1уд & \centering 0,001...100,000 & \centering Ом/км & \centering \arraybackslash 0,001 \\
\hline
\centering Удельное активное сопротивление нулевой последовательности ()  & \centering R0уд & \centering 0,001...100,000 & \centering Ом/км & \centering \arraybackslash 0,001 \\
\hline
\centering Удельное реактивное сопротивление нулевой последовательности ()  & \centering X0уд & \centering 0,001...100,000 & \centering Ом/км & \centering \arraybackslash 0,001 \\
\hline
%===t1
\end{longtable}
\normalsize


%%%%%%%%%%%%%%%%%%%%%%%%%%%%%%%%%%%%%%%%%%%%%%%%%%%%%%%%%% БК  %%%%%%%%%%%%%%%%%%%%%%%%%%%%%%%%%%%%%%%%%%%%%%%%%%%%%%%%%%%%%%%%%%
\par
В таблице \ref{bk:tbl1} приведены параметры, необходимые для настройки БК.

\footnotesize
\begin{longtable}{|>{\centering\arraybackslash}m{5.3cm}|>{\centering\arraybackslash}m{3.3cm}|>{\centering\arraybackslash}m{4.2cm}|>{\centering\arraybackslash}m{1.8cm}|>{\centering\arraybackslash}m{1cm}|}
\caption{Параметры для настройки БК\hfill\vspace{-0.5\baselineskip}}\label{bk:tbl1}\\ 
\hline
\rowcolor{gray!20}
Параметр (Параметр на ИЧМ) & Условное обозначение на схеме & Значение/ Диапазон & Единица измерения & Шаг \\ 
\hline
\endfirsthead
\caption*{\hspace{3pt}\emph{Продолжение таблицы \ref{bk:tbl1}\hfill\vspace{-0.5\baselineskip}}} \\ % сделано по ГОСТ 2.105 п.6.8.7
\hline
\rowcolor{gray!20}
Параметр (Параметр на ИЧМ) & Условное обозначение на схеме & Значение/ Диапазон & Единица измерения & Шаг \\ 
%\hline
\endhead
%\hline
\endfoot
%\hline
\endlastfoot
%==+t1*RDIS1|LVPSB
\centering Ввод функции в работу (Ввод\_функции)  & \centering SGF1 & \centering Не предусмотрено \\ Предусмотрено & \centering -- & \centering \arraybackslash -- \\
\hline
\centering Приращение тока ПП чувствительного органа (dI1\_чув)  & \centering dI1/dt> & \centering 0,08...3,00 & \centering о.е. & \centering \arraybackslash 0,01 \\
\hline
\centering Приращение тока ПП грубого органа (dI1\_груб)  & \centering dI1/dt\verb|>>| & \centering 0,12...5,00 & \centering о.е. & \centering \arraybackslash 0,01 \\
\hline
\centering Приращение тока ОП чувствительного органа (dI2\_чув)  & \centering dI2/dt> & \centering 0,04...1,50 & \centering о.е. & \centering \arraybackslash 0,01 \\
\hline
\centering Приращение тока ОП грубого органа (dI2\_груб)  & \centering dI2/dt\verb|>>| & \centering 0,06...2,50 & \centering о.е. & \centering \arraybackslash 0,01 \\
\hline
\centering Время ввода быстродействующих ступеней ДЗ от чувствительных органов (Тб\_чув)  & \centering T1 & \centering 200...1000 & \centering мс & \centering \arraybackslash 10 \\
\hline
\centering Время ввода быстродействующих ступеней ДЗ от грубых органов (Тб\_груб)  & \centering T2 & \centering 200...1000 & \centering мс & \centering \arraybackslash 10 \\
\hline
\centering Время ввода медленнодействующих ступеней ДЗ (Тм)  & \centering T3 & \centering 1000...15000 & \centering мс & \centering \arraybackslash 10 \\
\hline
\centering Ускоренный возврат логики БК при отключении выключателя (Ускор\_возврат)  & \centering SGF2 & \centering Не предусмотрено \\ Предусмотрено & \centering -- & \centering \arraybackslash -- \\
\hline
%===t1
\hline % ВНИМАНИЕ !!! Эта строка добавлена из-за висячей строки таблицы без заголовка в РЭ М300-ВВ 14-01-26, если будет линия на новой странице от таблицы - эту строку нужно закомметнировать
\end{longtable}
\normalsize


%%%%%%%%%%%%%%%%%%%%%%%%%%%%%%%%%%%%%%%%%%%%%%%%%%%%%%%%%% ТО %%%%%%%%%%%%%%%%%%%%%%%%%%%%%%%%%%%%%%%%%%%%%%%%%%%%%%%%%%%%%%%%%%
\par
В таблице \ref{to:tbl1} приведены параметры, необходимые для настройки ТО.

\footnotesize
\begin{longtable}{|>{\centering\arraybackslash}m{5cm}|>{\centering\arraybackslash}m{4cm}|>{\centering\arraybackslash}m{4cm}|>{\centering\arraybackslash}m{2cm}|>{\centering\arraybackslash}m{1cm}|}
\caption{Параметры для настройки ТО\hfill\vspace{-0.5\baselineskip}}\label{to:tbl1}\\ 
\hline
\rowcolor{gray!20}
Параметр (Параметр на ИЧМ) & Условное обозначение на схеме & Значение/ Диапазон & Единица измерения & Шаг \\ 
\hline
\endfirsthead
\caption*{\hspace{3pt}\emph{Продолжение таблицы \ref{to:tbl1}\hfill\vspace{-0.5\baselineskip}}} \\ % сделано по ГОСТ 2.105 п.6.8.7
\hline
\rowcolor{gray!20}
Параметр (Параметр на ИЧМ) & Условное обозначение на схеме & Значение/ Диапазон & Единица измерения & Шаг \\ 
%\hline
\endhead
%\hline
\endfoot
%\hline
\endlastfoot
%==+t1*PTOC1|LVTOC
\centering Ввод функции в работу (Ввод\_функции)  & \centering SGF1 & \centering Не предусмотрено \\ Предусмотрено & \centering -- & \centering \arraybackslash -- \\
\hline
\centering Ток срабатывания (Iср)  & \centering I> & \centering 0,10...30,00 & \centering о.е. & \centering \arraybackslash 0,01 \\
\hline
\centering Выдержка времени срабатывания (Tср)  & \centering T1 & \centering 0...10000 & \centering мс & \centering \arraybackslash 10 \\
\hline
%===t1
\end{longtable}
\normalsize


%%%%%%%%%%%%%%%%%%%%%%%%%%%%%%%%%%%%%%%%%%%%%%%%%%%%%%%%%% МТЗ %%%%%%%%%%%%%%%%%%%%%%%%%%%%%%%%%%%%%%%%%%%%%%%%%%%%%%%%%%%%%%%%%%
\par
В таблице \ref{mtz:tbl1} приведены параметры, необходимые для настройки МТЗ.

\footnotesize
\begin{longtable}{|>{\centering\arraybackslash}m{5.3cm}|>{\centering\arraybackslash}m{3.3cm}|>{\centering\arraybackslash}m{4.2cm}|>{\centering\arraybackslash}m{1.8cm}|>{\centering\arraybackslash}m{1cm}|}
\caption{Параметры для настройки ступеней МТЗ\hfill\vspace{-0.5\baselineskip}}\label{mtz:tbl1}\\ 
\hline
\rowcolor{gray!20}
Параметр (Параметр на ИЧМ) & Условное обозначение на схеме & Значение/ Диапазон & Единица измерения & Шаг \\ 
\hline
\endfirsthead
\caption*{\hspace{3pt}\emph{Продолжение таблицы \ref{mtz:tbl1}\hfill\vspace{-0.5\baselineskip}}} \\ % сделано по ГОСТ 2.105 п.6.8.7
\hline
\rowcolor{gray!20}
Параметр (Параметр на ИЧМ) & Условное обозначение на схеме & Значение/ Диапазон & Единица измерения & Шаг \\ 
%\hline
\endhead
%\hline
\endfoot
%\hline
\endlastfoot
%==+t1*PTOC1|LVTUVTOC
\multicolumn{5}{|c|}{ МТЗ 1 ст } \\ \hline 
\centering Ввод функции в работу (Ввод\_функции)  & \centering SGF1 & \centering Не предусмотрено \\ Предусмотрено & \centering -- & \centering \arraybackslash -- \\
\hline
\centering Ток срабатывания (Iср)  & \centering I> & \centering 0,10...30,00 & \centering о.е. & \centering \arraybackslash 0,01 \\
\hline
\centering Ток срабатывания грубого органа в режиме управляющего напряжения (Iср\_груб.)  & \centering -- & \centering 0,10...30,00 & \centering о.е. & \centering \arraybackslash 0,01 \\
\hline
\centering Ток срабатывания функции включения под нагрузку (Iфвн)  & \centering -- & \centering 0,10...30,00 & \centering о.е. & \centering \arraybackslash 0,01 \\
\hline
\centering Ввод функции включения на "холодную" нагрузку (Ввод\_ФВН)  & \centering SGF2 & \centering Не предусмотрено \\ Предусмотрено & \centering -- & \centering \arraybackslash -- \\
\hline
\centering Время отключенного состояния выключателя (Тзагруб.)  & \centering T2 & \centering 0...20000 & \centering мс & \centering \arraybackslash 10 \\
\hline
\centering Время ввода загрубления (Тсброс.загруб.)  & \centering T3 & \centering 0...20000 & \centering мс & \centering \arraybackslash 10 \\
\hline
\centering Тип пуска по напряжению (Тип\_КПОН)  & \centering SGF6 & \centering Управляющее напряжение \\ Вольтметровая блокировка & \centering -- & \centering \arraybackslash -- \\
\hline
\centering Режим контроля от БНТ (Реж\_БНТ)  & \centering SGF7 & \centering Не предусмотрено \\ Предусмотрено & \centering -- & \centering \arraybackslash -- \\
\hline
\centering Режим направленности (Направленность)  & \centering SGF3 & \centering Ненаправленная \\ Прямонаправленная \\ Обратнонаправленная & \centering -- & \centering \arraybackslash -- \\
\hline
\centering Режим ОНМ при неисправности ЦН (Реж\_БНН\_ОНМ)  & \centering SGF4 & \centering Блокировка \\ Деблокировка & \centering -- & \centering \arraybackslash -- \\
\hline
\centering Вывод направленности при АУ (Ненапр\_АУ)  & \centering SGF5 & \centering Не предусмотрено \\ Предусмотрено & \centering -- & \centering \arraybackslash -- \\
\hline
\centering Режим КПОН при неисправности ЦН (Реж\_БНН\_КПОН)  & \centering SGF8 & \centering Деблокировка \\ Блокировка & \centering -- & \centering \arraybackslash -- \\
\hline
\centering Режим контроля от КПОН (Реж\_КПОН)  & \centering SGF9 & \centering Не предусмотрено \\ Предусмотрено & \centering -- & \centering \arraybackslash -- \\
\hline
\centering Характеристика срабатывания (Характеристика)  & \centering -- & \centering Кривая F \\ Кривая E \\ Кривая D \\ Кривая B \\ Кривая A \\ Кривая C \\ Независимая \\ Кривая RXIDG & \centering -- & \centering \arraybackslash -- \\
\hline
\centering Коэффициент регулировки времени срабатывания зависимой характеристики (К)  & \centering -- & \centering 0,05...15,00 & \centering -- & \centering \arraybackslash 0,01 \\
\hline
\centering Независимая выдержка времени срабатывания (Tср)  & \centering T1 & \centering 0...100000 & \centering мс & \centering \arraybackslash 10 \\
\hline
\multicolumn{5}{|c|}{ МТЗ 2 ст } \\ \hline 
\centering Ввод функции в работу (Ввод\_функции)  & \centering SGF1 & \centering Не предусмотрено \\ Предусмотрено & \centering -- & \centering \arraybackslash -- \\
\hline
\centering Ток срабатывания (Iср)  & \centering I> & \centering 0,10...30,00 & \centering о.е. & \centering \arraybackslash 0,01 \\
\hline
\centering Ток срабатывания грубого органа в режиме управляющего напряжения (Iср\_груб.)  & \centering -- & \centering 0,10...30,00 & \centering о.е. & \centering \arraybackslash 0,01 \\
\hline
\centering Ток срабатывания функции включения под нагрузку (Iфвн)  & \centering -- & \centering 0,10...30,00 & \centering о.е. & \centering \arraybackslash 0,01 \\
\hline
\centering Ввод функции включения на "холодную" нагрузку (Ввод\_ФВН)  & \centering SGF2 & \centering Не предусмотрено \\ Предусмотрено & \centering -- & \centering \arraybackslash -- \\
\hline
\centering Время отключенного состояния выключателя (Тзагруб.)  & \centering T2 & \centering 0...20000 & \centering мс & \centering \arraybackslash 10 \\
\hline
\centering Время ввода загрубления (Тсброс.загруб.)  & \centering T3 & \centering 0...20000 & \centering мс & \centering \arraybackslash 10 \\
\hline
\centering Тип пуска по напряжению (Тип\_КПОН)  & \centering SGF6 & \centering Управляющее напряжение \\ Вольтметровая блокировка & \centering -- & \centering \arraybackslash -- \\
\hline
\centering Режим контроля от БНТ (Реж\_БНТ)  & \centering SGF7 & \centering Не предусмотрено \\ Предусмотрено & \centering -- & \centering \arraybackslash -- \\
\hline
\centering Режим направленности (Направленность)  & \centering SGF3 & \centering Ненаправленная \\ Прямонаправленная \\ Обратнонаправленная & \centering -- & \centering \arraybackslash -- \\
\hline
\centering Режим ОНМ при неисправности ЦН (Реж\_БНН\_ОНМ)  & \centering SGF4 & \centering Блокировка \\ Деблокировка & \centering -- & \centering \arraybackslash -- \\
\hline
\centering Вывод направленности при АУ (Ненапр\_АУ)  & \centering SGF5 & \centering Не предусмотрено \\ Предусмотрено & \centering -- & \centering \arraybackslash -- \\
\hline
\centering Режим КПОН при неисправности ЦН (Реж\_БНН\_КПОН)  & \centering SGF8 & \centering Деблокировка \\ Блокировка & \centering -- & \centering \arraybackslash -- \\
\hline
\centering Режим контроля от КПОН (Реж\_КПОН)  & \centering SGF9 & \centering Не предусмотрено \\ Предусмотрено & \centering -- & \centering \arraybackslash -- \\
\hline
\centering Характеристика срабатывания (Характеристика)  & \centering -- & \centering Кривая F \\ Кривая E \\ Кривая D \\ Кривая B \\ Кривая A \\ Кривая C \\ Независимая \\ Кривая RXIDG & \centering -- & \centering \arraybackslash -- \\
\hline
\centering Коэффициент регулировки времени срабатывания зависимой характеристики (К)  & \centering -- & \centering 0,05...15,00 & \centering -- & \centering \arraybackslash 0,01 \\
\hline
\centering Независимая выдержка времени срабатывания (Tср)  & \centering T1 & \centering 0...100000 & \centering мс & \centering \arraybackslash 10 \\
\hline
\multicolumn{5}{|c|}{ МТЗ 3 ст } \\ \hline 
\centering Ввод функции в работу (Ввод\_функции)  & \centering SGF1 & \centering Не предусмотрено \\ Предусмотрено & \centering -- & \centering \arraybackslash -- \\
\hline
\centering Ток срабатывания (Iср)  & \centering I> & \centering 0,10...30,00 & \centering о.е. & \centering \arraybackslash 0,01 \\
\hline
\centering Ток срабатывания грубого органа в режиме управляющего напряжения (Iср\_груб.)  & \centering -- & \centering 0,10...30,00 & \centering о.е. & \centering \arraybackslash 0,01 \\
\hline
\centering Ток срабатывания функции включения под нагрузку (Iфвн)  & \centering -- & \centering 0,10...30,00 & \centering о.е. & \centering \arraybackslash 0,01 \\
\hline
\centering Ввод функции включения на "холодную" нагрузку (Ввод\_ФВН)  & \centering SGF2 & \centering Не предусмотрено \\ Предусмотрено & \centering -- & \centering \arraybackslash -- \\
\hline
\centering Время отключенного состояния выключателя (Тзагруб.)  & \centering T2 & \centering 0...20000 & \centering мс & \centering \arraybackslash 10 \\
\hline
\centering Время ввода загрубления (Тсброс.загруб.)  & \centering T3 & \centering 0...20000 & \centering мс & \centering \arraybackslash 10 \\
\hline
\centering Тип пуска по напряжению (Тип\_КПОН)  & \centering SGF6 & \centering Управляющее напряжение \\ Вольтметровая блокировка & \centering -- & \centering \arraybackslash -- \\
\hline
\centering Режим контроля от БНТ (Реж\_БНТ)  & \centering SGF7 & \centering Не предусмотрено \\ Предусмотрено & \centering -- & \centering \arraybackslash -- \\
\hline
\centering Режим направленности (Направленность)  & \centering SGF3 & \centering Ненаправленная \\ Прямонаправленная \\ Обратнонаправленная & \centering -- & \centering \arraybackslash -- \\
\hline
\centering Режим ОНМ при неисправности ЦН (Реж\_БНН\_ОНМ)  & \centering SGF4 & \centering Блокировка \\ Деблокировка & \centering -- & \centering \arraybackslash -- \\
\hline
\centering Вывод направленности при АУ (Ненапр\_АУ)  & \centering SGF5 & \centering Не предусмотрено \\ Предусмотрено & \centering -- & \centering \arraybackslash -- \\
\hline
\centering Режим КПОН при неисправности ЦН (Реж\_БНН\_КПОН)  & \centering SGF8 & \centering Деблокировка \\ Блокировка & \centering -- & \centering \arraybackslash -- \\
\hline
\centering Режим контроля от КПОН (Реж\_КПОН)  & \centering SGF9 & \centering Не предусмотрено \\ Предусмотрено & \centering -- & \centering \arraybackslash -- \\
\hline
\centering Характеристика срабатывания (Характеристика)  & \centering -- & \centering Кривая F \\ Кривая E \\ Кривая D \\ Кривая B \\ Кривая A \\ Кривая C \\ Независимая \\ Кривая RXIDG & \centering -- & \centering \arraybackslash -- \\
\hline
\centering Коэффициент регулировки времени срабатывания зависимой характеристики (К)  & \centering -- & \centering 0,05...15,00 & \centering -- & \centering \arraybackslash 0,01 \\
\hline
\centering Независимая выдержка времени срабатывания (Tср)  & \centering T1 & \centering 0...100000 & \centering мс & \centering \arraybackslash 10 \\
\hline
\multicolumn{5}{|c|}{ КПОН } \\ \hline 
\centering Напряжение срабатывания (Uср)  & \centering U< & \centering 2,0...100,0 & \centering В & \centering \arraybackslash 0,1 \\
\hline
\centering Напряжение срабатывания ОП (U2ср)  & \centering U2> & \centering 2,0...50,0 & \centering В & \centering \arraybackslash 0,1 \\
\hline
\centering Режим пуска (Реж\_пуска)  & \centering SGF1 & \centering По Uмин \\ Комбинированный \\ Внешний & \centering -- & \centering \arraybackslash -- \\
\hline
\multicolumn{5}{|c|}{ ОНМ } \\ \hline 
\centering Угол максимальной чувствительности (Фмч)  & \centering φмч & \centering -179,0...180,0 & \centering град. & \centering \arraybackslash 0,1 \\
\hline
\multicolumn{5}{|c|}{ БНТ } \\ \hline 
\centering Отношение тока второй гармоники к току основной гармоники (Ih2/Ih1)  & \centering I2г/I1г> & \centering 5...100 & \centering \% & \centering \arraybackslash 1 \\
\hline
\centering Ток блокировки (Iср)  & \centering I> & \centering 0,1...15,0 & \centering о.е. & \centering \arraybackslash 0,1 \\
\hline
\centering Перекрестная блокировка (Перекрест\_блок)  & \centering SGF1 & \centering Не предусмотрено \\ Предусмотрено & \centering -- & \centering \arraybackslash -- \\
\hline
\multicolumn{5}{|c|}{ БЛЗ } \\ \hline 
\centering Выбор ступени блокировки (БлокЛЗ\_выбор\_ст)  & \centering SGF1 & \centering Не предусмотрено \\ 1 ступень \\ 2 ступень \\ 3 ступень & \centering -- & \centering \arraybackslash -- \\
\hline
\multicolumn{5}{|c|}{ АУ МТЗ } \\ \hline 
\centering Выбор ускоряемой ступени (Выбор\_ст)  & \centering SGF1 & \centering 1 ступень \\ 2 ступень \\ 3 ступень & \centering -- & \centering \arraybackslash -- \\
\hline
\multicolumn{5}{|c|}{ ОУ МТЗ } \\ \hline 
\centering Выбор оперативно ускоряемой ступени (Выбор\_ст\_ОУ)  & \centering SGF1 & \centering 2 ступень \\ 3 ступень & \centering -- & \centering \arraybackslash -- \\
\hline
\centering Выдержка времени срабатывания оперативного ускорения (T\_ОУ)  & \centering T1 & \centering 0...3000 & \centering мс & \centering \arraybackslash 10 \\
\hline
%===t1
\end{longtable}
\normalsize


%%%%%%%%%%%%%%%%%%%%%%%%%%%%%%%%%%%%%%%%%%%%%%%%%%%%%%%%%% АУ %%%%%%%%%%%%%%%%%%%%%%%%%%%%%%%%%%%%%%%%%%%%%%%%%%%%%%%%%%%%%%%%%%
\par
В таблице \ref{au:tbl1} приведены параметры, необходимые для настройки АУ.

\footnotesize
\begin{longtable}{|>{\centering\arraybackslash}m{5.3cm}|>{\centering\arraybackslash}m{3.3cm}|>{\centering\arraybackslash}m{4.2cm}|>{\centering\arraybackslash}m{1.8cm}|>{\centering\arraybackslash}m{1cm}|}
\caption{Параметры для настройки АУ\hfill\vspace{-0.5\baselineskip}}\label{au:tbl1}\\ 
\hline
\rowcolor{gray!20}
Параметр (Параметр на ИЧМ) & Условное обозначение на схеме & Значение/ Диапазон & Единица измерения & Шаг \\ 
\hline
\endfirsthead
\caption*{\hspace{3pt}\emph{Продолжение таблицы \ref{au:tbl1}\hfill\vspace{-0.5\baselineskip}}} \\ % сделано по ГОСТ 2.105 п.6.8.7
\hline
\rowcolor{gray!20}
Параметр (Параметр на ИЧМ) & Условное обозначение на схеме & Значение/ Диапазон & Единица измерения & Шаг \\ 
%\hline
\endhead
%\hline
\endfoot
%\hline
\endlastfoot
%==+t1*PSOF1|LVLINAUA
\centering Ввод функции в работу (Ввод\_функции)  & \centering SGF1 & \centering Не предусмотрено \\ Предусмотрено & \centering -- & \centering \arraybackslash -- \\
\hline
\centering Автоматическое ускорение МТЗ (АУ\_МТЗ)  & \centering SGF2 & \centering Не предусмотрено \\ Предусмотрено & \centering -- & \centering \arraybackslash -- \\
\hline
\centering Автоматическое ускорение ДЗ (АУ\_ДЗ)  & \centering SGF3 & \centering Не предусмотрено \\ Предусмотрено & \centering -- & \centering \arraybackslash -- \\
\hline
\centering Автоматическое ускорение ТЗНП (АУ\_ТЗНП)  & \centering SGF4 & \centering Не предусмотрено \\ Предусмотрено & \centering -- & \centering \arraybackslash -- \\
\hline
\centering Выдержка времени срабатывания ускоряемой ступени МТЗ (Тср\_МТЗ)  & \centering T1 & \centering 0...3000 & \centering мс & \centering \arraybackslash 10 \\
\hline
\centering Выдержка времени срабатывания ускоряемой ступени ДЗ (Тср\_ДЗ)  & \centering T2 & \centering 0...3000 & \centering мс & \centering \arraybackslash 10 \\
\hline
\centering Выдержка времени срабатывания ускоряемой ступени ТЗНП (Тср\_ТЗНП)  & \centering T3 & \centering 0...3000 & \centering мс & \centering \arraybackslash 10 \\
\hline
\centering Время ввода ускорения при включении выключателя (Твв)  & \centering T4 & \centering 0...5000 & \centering мс & \centering \arraybackslash 10 \\
\hline
%===t1
\end{longtable}
\normalsize

%%%%%%%%%%%%%%%%%%%%%%%%%%%%%%%%%%%%%%%%%%%%%%%%%%%%%%%%%% ТЗНП %%%%%%%%%%%%%%%%%%%%%%%%%%%%%%%%%%%%%%%%%%%%%%%%%%%%%%%%%%%%%%%%%%
\par
В таблице \ref{tznp:tbl1} приведены параметры, необходимые для настройки ТЗНП.

\footnotesize
\begin{longtable}{|>{\centering\arraybackslash}m{5.3cm}|>{\centering\arraybackslash}m{3.3cm}|>{\centering\arraybackslash}m{4.2cm}|>{\centering\arraybackslash}m{1.8cm}|>{\centering\arraybackslash}m{1cm}|}
\caption{Параметры для настройки ступеней ТЗНП\hfill\vspace{-0.5\baselineskip}}\label{tznp:tbl1}\\ 
\hline
\rowcolor{gray!20}
Параметр (Параметр на ИЧМ) & Условное обозначение на схеме & Значение/ Диапазон & Единица измерения & Шаг \\ 
\hline
\endfirsthead
\caption*{\hspace{3pt}\emph{Продолжение таблицы \ref{tznp:tbl1}\hfill\vspace{-0.5\baselineskip}}} \\ % сделано по ГОСТ 2.105 п.6.8.7
\hline
\rowcolor{gray!20}
Параметр (Параметр на ИЧМ) & Условное обозначение на схеме & Значение/ Диапазон & Единица измерения & Шаг \\ 
%\hline
\endhead
%\hline
\endfoot
%\hline
\endlastfoot
%==+t1*ZSPTOC1|LVZSTOC
\multicolumn{5}{|c|}{ ТЗНП 1 ст } \\ \hline 
\centering Ввод функции в работу (Ввод\_функции)  & \centering SGF1 & \centering Не предусмотрено \\ Предусмотрено & \centering -- & \centering \arraybackslash -- \\
\hline
\centering Режим оперативного блокирования ЧС (Реж\_БЧС)  & \centering SGF8 & \centering Не предусмотрено \\ Предусмотрено & \centering -- & \centering \arraybackslash -- \\
\hline
\centering Режим контроля от БНТ (Реж\_БНТ)  & \centering SGF7 & \centering Не предусмотрено \\ Предусмотрено & \centering -- & \centering \arraybackslash -- \\
\hline
\centering Ток срабатывания НП (3I0ср)  & \centering 3I0> & \centering 0,10...30,00 & \centering о.е. & \centering \arraybackslash 0,01 \\
\hline
\centering Вывод направленности при АУ (Ненапр\_АУ)  & \centering SGF4 & \centering Не предусмотрено \\ Предусмотрено & \centering -- & \centering \arraybackslash -- \\
\hline
\centering Режим направленности (Направленность)  & \centering SGF2 & \centering Ненаправленная \\ Прямонаправленная \\ Обратнонаправленная & \centering -- & \centering \arraybackslash -- \\
\hline
\centering Режим ОНМ НП при неисправности ЦН (Реж\_БНН\_ОНМ\_НП)  & \centering SGF3 & \centering Блокировка \\ Деблокировка & \centering -- & \centering \arraybackslash -- \\
\hline
\centering Режим ПО 3U0 при неисправности ЦН (Реж\_БНН\_3U0)  & \centering SGF5 & \centering Блокировка \\ Деблокировка & \centering -- & \centering \arraybackslash -- \\
\hline
\centering Напряжение срабатывания НП (3U0ср)  & \centering 3U0> & \centering 2,0...100,0 & \centering В & \centering \arraybackslash 0,1 \\
\hline
\centering Режим контроля от ПО 3U0 (Реж\_3U0)  & \centering SGF6 & \centering Не предусмотрено \\ Предусмотрено & \centering -- & \centering \arraybackslash -- \\
\hline
\centering Характеристика срабатывания (Характеристика)  & \centering -- & \centering Кривая F \\ Кривая E \\ Кривая D \\ Кривая B \\ Кривая A \\ Кривая C \\ Независимая \\ Кривая RXIDG & \centering -- & \centering \arraybackslash -- \\
\hline
\centering Коэффициент регулировки времени срабатывания зависимой характеристики (К)  & \centering -- & \centering 0,05...15,00 & \centering -- & \centering \arraybackslash 0,01 \\
\hline
\centering Независимая выдержка времени срабатывания (Tср)  & \centering T1 & \centering 0...100000 & \centering мс & \centering \arraybackslash 10 \\
\hline
\multicolumn{5}{|c|}{ ТЗНП 2 ст } \\ \hline 
\centering Ввод функции в работу (Ввод\_функции)  & \centering SGF1 & \centering Не предусмотрено \\ Предусмотрено & \centering -- & \centering \arraybackslash -- \\
\hline
\centering Режим оперативного блокирования ЧС (Реж\_БЧС)  & \centering SGF8 & \centering Не предусмотрено \\ Предусмотрено & \centering -- & \centering \arraybackslash -- \\
\hline
\centering Режим контроля от БНТ (Реж\_БНТ)  & \centering SGF7 & \centering Не предусмотрено \\ Предусмотрено & \centering -- & \centering \arraybackslash -- \\
\hline
\centering Ток срабатывания НП (3I0ср)  & \centering 3I0> & \centering 0,10...30,00 & \centering о.е. & \centering \arraybackslash 0,01 \\
\hline
\centering Вывод направленности при АУ (Ненапр\_АУ)  & \centering SGF4 & \centering Не предусмотрено \\ Предусмотрено & \centering -- & \centering \arraybackslash -- \\
\hline
\centering Режим направленности (Направленность)  & \centering SGF2 & \centering Ненаправленная \\ Прямонаправленная \\ Обратнонаправленная & \centering -- & \centering \arraybackslash -- \\
\hline
\centering Режим ОНМ НП при неисправности ЦН (Реж\_БНН\_ОНМ\_НП)  & \centering SGF3 & \centering Блокировка \\ Деблокировка & \centering -- & \centering \arraybackslash -- \\
\hline
\centering Режим ПО 3U0 при неисправности ЦН (Реж\_БНН\_3U0)  & \centering SGF5 & \centering Блокировка \\ Деблокировка & \centering -- & \centering \arraybackslash -- \\
\hline
\centering Напряжение срабатывания НП (3U0ср)  & \centering 3U0> & \centering 2,0...100,0 & \centering В & \centering \arraybackslash 0,1 \\
\hline
\centering Режим контроля от ПО 3U0 (Реж\_3U0)  & \centering SGF6 & \centering Не предусмотрено \\ Предусмотрено & \centering -- & \centering \arraybackslash -- \\
\hline
\centering Характеристика срабатывания (Характеристика)  & \centering -- & \centering Кривая F \\ Кривая E \\ Кривая D \\ Кривая B \\ Кривая A \\ Кривая C \\ Независимая \\ Кривая RXIDG & \centering -- & \centering \arraybackslash -- \\
\hline
\centering Коэффициент регулировки времени срабатывания зависимой характеристики (К)  & \centering -- & \centering 0,05...15,00 & \centering -- & \centering \arraybackslash 0,01 \\
\hline
\centering Независимая выдержка времени срабатывания (Tср)  & \centering T1 & \centering 0...100000 & \centering мс & \centering \arraybackslash 10 \\
\hline
\multicolumn{5}{|c|}{ АУ ТЗНП } \\ \hline 
\centering Выбор ускоряемой ступени (Выбор\_ст)  & \centering SGF1 & \centering 1 ступень \\ 2 ступень & \centering -- & \centering \arraybackslash -- \\
\hline
\multicolumn{5}{|c|}{ ОНМ НП } \\ \hline 
\centering Угол максимальной чувствительности (Фмч)  & \centering φмч & \centering -179,0...180,0 & \centering град. & \centering \arraybackslash 0,1 \\
\hline
\multicolumn{5}{|c|}{ БНТ НП } \\ \hline 
\centering Отношение тока второй гармоники к току основной гармоники (Ih2/Ih1)  & \centering I2г/I1г> & \centering 5...100 & \centering \% & \centering \arraybackslash 1 \\
\hline
\centering Ток блокировки НП (3I0ср)  & \centering 3I0> & \centering 0,1...15,0 & \centering о.е. & \centering \arraybackslash 0,1 \\
\hline
\multicolumn{5}{|c|}{ ОУ ТЗНП } \\ \hline 
\centering Выбор оперативно ускоряемой ступени (Выбор\_ст\_ОУ)  & \centering SGF1 & \centering 1 ступень \\ 2 ступень & \centering -- & \centering \arraybackslash -- \\
\hline
\centering Выдержка времени срабатывания оперативного ускорения (T\_ОУ)  & \centering T1 & \centering 0...3000 & \centering мс & \centering \arraybackslash 10 \\
\hline
%===t1
\end{longtable}
\normalsize

%%%%%%%%%%%%%%%%%%%%%%%%%%%%%%%%%%%%%%%%%%%%%%%%%%%%%%%%%% ТЗНП %%%%%%%%%%%%%%%%%%%%%%%%%%%%%%%%%%%%%%%%%%%%%%%%%%%%%%%%%%%%%%%%%%
\par
В таблице \ref{nzozz:tbl1} приведены параметры, необходимые для настройки НЗОЗЗ.

\footnotesize
\begin{longtable}{|>{\centering\arraybackslash}m{5.3cm}|>{\centering\arraybackslash}m{3.3cm}|>{\centering\arraybackslash}m{4.2cm}|>{\centering\arraybackslash}m{1.8cm}|>{\centering\arraybackslash}m{1cm}|}
\caption{Параметры для настройки ПО 3I0\hfill\vspace{-0.5\baselineskip}}\label{nzozz:tbl1}\\ 
\hline
\rowcolor{gray!20}
Параметр (Параметр на ИЧМ) & Условное обозначение на схеме & Значение/ Диапазон & Единица измерения & Шаг \\ 
\hline
\endfirsthead
\caption*{\hspace{3pt}\emph{Продолжение таблицы \ref{nzozz:tbl1}\hfill\vspace{-0.5\baselineskip}}} \\ % сделано по ГОСТ 2.105 п.6.8.7
\hline
\rowcolor{gray!20}
Параметр (Параметр на ИЧМ) & Условное обозначение на схеме & Значение/ Диапазон & Единица измерения & Шаг \\ 
%\hline
\endhead
%\hline
\endfoot
%\hline
\endlastfoot
%==+t1*ZSPSDE1|LVTRSDE
\multicolumn{5}{|c|}{ ВыхЛогНЗОЗЗ } \\ \hline 
\centering Ввод функции в работу (Ввод\_функции)  & \centering SGF1 & \centering Не предусмотрено \\ Предусмотрено & \centering -- & \centering \arraybackslash -- \\
\hline
\centering Режим работы выходной логики НЗОЗЗ (Реж\_НЗОЗЗ)  & \centering SGF2 & \centering По 3I0 \\ По 3I0 и 3U0 \\ По 3I0вг \\ По 3I0вг и 3U0 \\ По 3U0 & \centering -- & \centering \arraybackslash -- \\
\hline
\centering Отключение от логики обнаружения переходных/ перемежающихся замыканий (Откл\_от\_ЛОППЗ)  & \centering SGF3 & \centering Не предусмотрено \\ Предусмотрено & \centering -- & \centering \arraybackslash -- \\
\hline
\centering Характеристика срабатывания (Характ\_ср)  & \centering -- & \centering Кривая F \\ Кривая E \\ Кривая D \\ Кривая B \\ Кривая A \\ Кривая C \\ Независимая \\ Кривая RXIDG & \centering -- & \centering \arraybackslash -- \\
\hline
\centering Коэффициент регулировки времени срабатывания зависимой характеристики (Kср)  & \centering -- & \centering 0,05...15,00 & \centering -- & \centering \arraybackslash 0,01 \\
\hline
\centering Независимая выдержка времени срабатывания (Tср)  & \centering T1 & \centering 0...100000 & \centering мс & \centering \arraybackslash 10 \\
\hline
\centering Характеристика срабатывания на сигнал (Характ\_сигн)  & \centering -- & \centering Кривая F \\ Кривая E \\ Кривая D \\ Кривая B \\ Кривая A \\ Кривая C \\ Независимая \\ Кривая RXIDG & \centering -- & \centering \arraybackslash -- \\
\hline
\centering Коэффициент регулировки времени срабатывания на сигнал
зависимой характеристики (Kсигн)  & \centering -- & \centering 0,05...15,00 & \centering -- & \centering \arraybackslash 0,01 \\
\hline
\centering Независимая выдержка времени срабатывания на сигнал (Tсигн)  & \centering T2 & \centering 0...100000 & \centering мс & \centering \arraybackslash 10 \\
\hline
\multicolumn{5}{|c|}{ ПО 3I0 } \\ \hline 
\centering Ввод функции в работу (Ввод\_функции)  & \centering SGF1 & \centering Не предусмотрено \\ Предусмотрено & \centering -- & \centering \arraybackslash -- \\
\hline
\centering Ток срабатывания НП (3I0ср)  & \centering 3I0> & \centering 0,01...3,00 & \centering А & \centering \arraybackslash 0,01 \\
\hline
\centering Режим направленности (Направленность)  & \centering SGF2 & \centering Ненаправленная \\ Прямонаправленная \\ Обратнонаправленная & \centering -- & \centering \arraybackslash -- \\
\hline
\centering Работа ПО 3I0 при неисправности ЦН (Реж\_БНН\_3U0)  & \centering SGF3 & \centering Не предусмотрено \\ Блокировка \\ Деблокировка & \centering -- & \centering \arraybackslash -- \\
\hline
\multicolumn{5}{|c|}{ ПО 3I0вг } \\ \hline 
\centering Ввод функции в работу (Ввод\_функции)  & \centering SGF1 & \centering Не предусмотрено \\ Предусмотрено & \centering -- & \centering \arraybackslash -- \\
\hline
\centering Ток срабатывания высших гармоник (3I0вг)  & \centering 3I0 вг > & \centering 0,001...1,000 & \centering А & \centering \arraybackslash 0,001 \\
\hline
\multicolumn{5}{|c|}{ ПО 3U0 } \\ \hline 
\centering Ввод функции в работу (Ввод\_функции)  & \centering SGF1 & \centering Не предусмотрено \\ Предусмотрено & \centering -- & \centering \arraybackslash -- \\
\hline
\centering Напряжение срабатывания НП (3U0ср)  & \centering 3U0> & \centering 2,0...100,0 & \centering В & \centering \arraybackslash 0,1 \\
\hline
\centering Режим ПО 3U0 при неисправности ЦН (Реж\_БНН\_3U0)  & \centering SGF2 & \centering Не предусмотрено \\ Блокировка \\ Деблокировка & \centering -- & \centering \arraybackslash -- \\
\hline
\multicolumn{5}{|c|}{ ОНМ НП НЗОЗЗ } \\ \hline 
\centering Угол максимальной чувствительности (Фмч)  & \centering φмч & \centering -179,0...180,0 & \centering град. & \centering \arraybackslash 0,1 \\
\hline
\multicolumn{5}{|c|}{ ЛОППЗ } \\ \hline 
\centering Режим работы (Режим\_работы)  & \centering SGF1 & \centering Не предусмотрено \\ Предусмотрено & \centering -- & \centering \arraybackslash -- \\
\hline
\centering Ток срабатывания НП (3I0ср)  & \centering 3I0> & \centering 0,05...10,00 & \centering А & \centering \arraybackslash 0,01 \\
\hline
\centering Напряжение срабатывания НП (3U0ср)  & \centering 3U0> & \centering 10,0...150,0 & \centering В & \centering \arraybackslash 0,1 \\
\hline
\centering Направление замыкания (Направление)  & \centering SGF2 & \centering Прямое \\ Обратное & \centering -- & \centering \arraybackslash -- \\
\hline
\centering Число импульсов для пуска (Nср)  & \centering -- & \centering 2...100 & \centering -- & \centering \arraybackslash 1 \\
\hline
\centering Время отсутствия мощности НП (Tвоз)  & \centering T2 & \centering 1...1000 & \centering мс & \centering \arraybackslash 1 \\
\hline
\centering Независимая выдержка времени срабатывания (Tср)  & \centering T1 & \centering 0...100000 & \centering мс & \centering \arraybackslash 10 \\
\hline
%===t1
\end{longtable}
\normalsize


%%%%%%%%%%%%%%%%%%%%%%%%%%%%%%%%%%%%%%%%%%%%%%%%%%%%%%%%%% ЗДЗ %%%%%%%%%%%%%%%%%%%%%%%%%%%%%%%%%%%%%%%%%%%%%%%%%%%%%%%%%%%%%%%%%%
\par
В таблице \ref{zdz:tbl1} приведены параметры, необходимые для настройки функции.

\footnotesize
\begin{longtable}{|>{\centering\arraybackslash}m{5.3cm}|>{\centering\arraybackslash}m{3.3cm}|>{\centering\arraybackslash}m{4.2cm}|>{\centering\arraybackslash}m{1.8cm}|>{\centering\arraybackslash}m{1cm}|}
\caption{Параметры для настройки ЗДЗ\hfill\vspace{-0.5\baselineskip}}\label{zdz:tbl1}\\ 
\hline
\rowcolor{gray!20}
Параметр (Параметр на ИЧМ) & Условное обозначение на схеме & Значение/ Диапазон & Единица измерения & Шаг \\ 
\hline
\endfirsthead
\caption*{\hspace{3pt}\emph{Продолжение таблицы \ref{zdz:tbl1}\hfill\vspace{-0.5\baselineskip}}} \\ % сделано по ГОСТ 2.105 п.6.8.7
\hline
\rowcolor{gray!20}
Параметр (Параметр на ИЧМ) & Условное обозначение на схеме & Значение/ Диапазон & Единица измерения & Шаг \\ 
%\hline
\endhead
%\hline
\endfoot
%\hline
\endlastfoot
%==+t1*ARCPTOC1|ARCTOC
\centering Ввод функции в работу (Ввод\_функции)  & \centering SGF1 & \centering Не предусмотрено \\ Предусмотрено & \centering -- & \centering \arraybackslash -- \\
\hline
\centering Ток срабатывания (Iср)  & \centering I> & \centering 0,10...30,00 & \centering о.е. & \centering \arraybackslash 0,01 \\
\hline
\centering Напряжение срабатывания (3U0ср)  & \centering 3U0> & \centering 2,0...150,0 & \centering В & \centering \arraybackslash 0,1 \\
\hline
\centering Блокировка при неисправности ЗДЗ (Блок\_неиспр\_ЗДЗ)  & \centering SGF2 & \centering Не предусмотрено \\ Предусмотрено & \centering -- & \centering \arraybackslash -- \\
\hline
\centering Контроль тока (Контр\_тока)  & \centering SGF3 & \centering Не предусмотрено \\ Контроль I \\ Контроль I или 3U0 \\ Внешний & \centering -- & \centering \arraybackslash -- \\
\hline
\centering Режим контроля от БНН (Реж\_БНН)  & \centering SGF4 & \centering Без контроля \\ С контролем & \centering -- & \centering \arraybackslash -- \\
\hline
\centering Выдержка времени неисправности ЗДЗ (Tнеисп)  & \centering T2 & \centering 0...10000 & \centering мс & \centering \arraybackslash 10 \\
\hline
\centering Выдержка времени срабатывания (Tср)  & \centering T1 & \centering 0...10000 & \centering мс & \centering \arraybackslash 10 \\
\hline
%===t1
\end{longtable}
\normalsize


%%%%%%%%%%%%%%%%%%%%%%%%%%%%%%%%%%%%%%%%%%%%%%%%%%%%%%%%%% ЗОП %%%%%%%%%%%%%%%%%%%%%%%%%%%%%%%%%%%%%%%%%%%%%%%%%%%%%%%%%%%%%%%%%%
\par
В таблице \ref{zop:tbl1} приведены параметры, необходимые для настройки ЗОП.

\footnotesize
\begin{longtable}{|>{\centering\arraybackslash}m{5.3cm}|>{\centering\arraybackslash}m{3.3cm}|>{\centering\arraybackslash}m{4.2cm}|>{\centering\arraybackslash}m{1.8cm}|>{\centering\arraybackslash}m{1cm}|}
\caption{Параметры для настройки ЗОП\hfill\vspace{-0.5\baselineskip}}\label{zop:tbl1}\\ 
\hline
\rowcolor{gray!20}
Параметр (Параметр на ИЧМ) & Условное обозначение на схеме & Значение/ Диапазон & Единица измерения & Шаг \\ 
\hline
\endfirsthead
\caption*{\hspace{3pt}\emph{Продолжение таблицы \ref{zop:tbl1}\hfill\vspace{-0.5\baselineskip}}} \\ % сделано по ГОСТ 2.105 п.6.8.7
\hline
\rowcolor{gray!20}
Параметр (Параметр на ИЧМ) & Условное обозначение на схеме & Значение/ Диапазон & Единица измерения & Шаг \\ 
%\hline
\endhead
%\hline
\endfoot
%\hline
\endlastfoot
%==+t1*NSPTOC1|LVNSTOC
\centering Ввод функции в работу (Ввод\_функции)  & \centering SGF1 & \centering Не предусмотрено \\ Предусмотрено & \centering -- & \centering \arraybackslash -- \\
\hline
\centering Ток срабатывания ОП (I2ср)  & \centering I2> & \centering 0,04...4,00 & \centering о.е. & \centering \arraybackslash 0,01 \\
\hline
\centering Отношение тока ОП к току ПП (Kнесим)  & \centering I2/I1> & \centering 0,02...1,00 & \centering о.е. & \centering \arraybackslash 0,01 \\
\hline
\centering Режим работы ЗОП (Реж\_ЗОП)  & \centering SGF2 & \centering по I2 \\ по I2/I1 \\ по I2 или по I2/I1 & \centering -- & \centering \arraybackslash -- \\
\hline
\centering Выдержка времени срабатывания (Tср)  & \centering T1 & \centering 100...20000 & \centering мс & \centering \arraybackslash 10 \\
\hline
%===t1
\end{longtable}
\normalsize


%%%%%%%%%%%%%%%%%%%%%%%%%%%%%%%%%%%%%%%%%%%%%%%%%%%%%%%%%% ЗП %%%%%%%%%%%%%%%%%%%%%%%%%%%%%%%%%%%%%%%%%%%%%%%%%%%%%%%%%%%%%%%%%%
\par
В таблице \ref{zp:tbl1} приведены параметры, необходимые для настройки ЗП.

\footnotesize
\begin{longtable}{|>{\centering\arraybackslash}m{5.3cm}|>{\centering\arraybackslash}m{3.3cm}|>{\centering\arraybackslash}m{4.2cm}|>{\centering\arraybackslash}m{1.8cm}|>{\centering\arraybackslash}m{1cm}|}
\caption{Параметры для настройки ЗП\hfill\vspace{-0.5\baselineskip}}\label{zp:tbl1}\\ 
\hline
\rowcolor{gray!20}
Параметр (Параметр на ИЧМ) & Условное обозначение на схеме & Значение/ Диапазон & Единица измерения & Шаг \\ 
\hline
\endfirsthead
\caption*{\hspace{3pt}\emph{Продолжение таблицы \ref{zp:tbl1}\hfill\vspace{-0.5\baselineskip}}} \\ % сделано по ГОСТ 2.105 п.6.8.7
\hline
\rowcolor{gray!20}
Параметр (Параметр на ИЧМ) & Условное обозначение на схеме & Значение/ Диапазон & Единица измерения & Шаг \\ 
%\hline
\endhead
%\hline
\endfoot
%\hline
\endlastfoot
%==+t1*OVCPTOC1|AUTOVC
\centering Ввод функции в работу (Ввод\_функции)  & \centering SGF1 & \centering Не предусмотрено \\ Предусмотрено & \centering -- & \centering \arraybackslash -- \\
\hline
\centering Ток срабатывания (Iср)  & \centering I> & \centering 0,10...5,00 & \centering о.е. & \centering \arraybackslash 0,01 \\
\hline
\centering Выдержка времени срабатывания (Tср)  & \centering T1 & \centering 0...600000 & \centering мс & \centering \arraybackslash 100 \\
\hline
%===t1
\end{longtable}
\normalsize


%%%%%%%%%%%%%%%%%%%%%%%%%%%%%%%%%%%%%%%%%%%%%%%%%%%%%%%%%% ГСОЗЗ %%%%%%%%%%%%%%%%%%%%%%%%%%%%%%%%%%%%%%%%%%%%%%%%%%%%%%%%%%%%%%%%%%
\par
В таблице \ref{gsozz:tbl1} приведены параметры, необходимые для настройки ГСОЗЗ.

\footnotesize
\begin{longtable}{|>{\centering\arraybackslash}m{5.3cm}|>{\centering\arraybackslash}m{3.3cm}|>{\centering\arraybackslash}m{4.2cm}|>{\centering\arraybackslash}m{1.8cm}|>{\centering\arraybackslash}m{1cm}|}
\caption{Параметры для настройки ГСОЗЗ\hfill\vspace{-0.5\baselineskip}}\label{gsozz:tbl1}\\ 
\hline
\rowcolor{gray!20}
Параметр (Параметр на ИЧМ) & Условное обозначение на схеме & Значение/ Диапазон & Единица измерения & Шаг \\ 
\hline
\endfirsthead
\caption*{\hspace{3pt}\emph{Продолжение таблицы \ref{gsozz:tbl1}\hfill\vspace{-0.5\baselineskip}}} \\ % сделано по ГОСТ 2.105 п.6.8.7
\hline
\rowcolor{gray!20}
Параметр (Параметр на ИЧМ) & Условное обозначение на схеме & Значение/ Диапазон & Единица измерения & Шаг \\ 
%\hline
\endhead
%\hline
\endfoot
%\hline
\endlastfoot
%==+t1*PSDE1|LVSDE
\centering Ввод функции в работу (Ввод\_функции)  & \centering SGF1 & \centering Не предусмотрено \\ Предусмотрено & \centering -- & \centering \arraybackslash -- \\
\hline
\centering Ток срабатывания НП (3I0ср)  & \centering 3I0> & \centering 0,01...10,00 & \centering А & \centering \arraybackslash 0,01 \\
\hline
\centering Ток срабатывания высших гармоник (3I0вг)  & \centering 3I0 в.г.> & \centering 0,001...1,000 & \centering А & \centering \arraybackslash 0,001 \\
\hline
\centering Напряжение срабатывания НП (3U0ср)  & \centering 3U0> & \centering 2,0...100,0 & \centering В & \centering \arraybackslash 0,1 \\
\hline
\centering Режим работы ГСОЗЗ (Реж\_ГСОЗЗ)  & \centering SGF4 & \centering По 3I0 \\ По 3I0вг & \centering -- & \centering \arraybackslash -- \\
\hline
\centering Режим ПО 3U0 при неисправности ЦН (Реж\_при\_БНН)  & \centering SGF3 & \centering Блокировка \\ Деблокировка & \centering -- & \centering \arraybackslash -- \\
\hline
\centering Режим работы от ПО 3U0 (Работа\_от\_3U0)  & \centering SGF2 & \centering Не предусмотрено \\ Предусмотрено & \centering -- & \centering \arraybackslash -- \\
\hline
\centering Характеристика срабатывания (Характеристика)  & \centering -- & \centering Кривая F \\ Кривая E \\ Кривая D \\ Кривая B \\ Кривая A \\ Кривая C \\ Независимая \\ Кривая RXIDG & \centering -- & \centering \arraybackslash -- \\
\hline
\centering Коэффициент регулировки времени срабатывания зависимой характеристики (К)  & \centering -- & \centering 0,05...15,00 & \centering -- & \centering \arraybackslash 0,01 \\
\hline
\centering Независимая выдержка времени срабатывания (Tср)  & \centering T1 & \centering 0...100000 & \centering мс & \centering \arraybackslash 10 \\
\hline
%===t1
\end{longtable}
\normalsize



%%%%%%%%%%%%%%%%%%%%%%%%%%%%%%%%%%%%%%%%%%%%%%%%%%%%%%%%%% ЛО ГЗсигн %%%%%%%%%%%%%%%%%%%%%%%%%%%%%%%%%%%%%%%%%%%%%%%%%%%%%%%%%%%%%%%%%%
\par

В таблице \ref{logzsign:tbl1} приведены параметры, необходимые для настройки ЛО ГЗсигн.

\footnotesize
\begin{longtable}{|>{\centering\arraybackslash}m{5.3cm}|>{\centering\arraybackslash}m{3.3cm}|>{\centering\arraybackslash}m{4.2cm}|>{\centering\arraybackslash}m{1.8cm}|>{\centering\arraybackslash}m{1cm}|}
\caption{Параметры для настройки логики отключения при срабатывании сигнального контакта газового реле трансформатора\hfill\vspace{-0.5\baselineskip}}\label{logzsign:tbl1}\\ 
\hline
\rowcolor{gray!20}
Параметр (Параметр на ИЧМ) & Условное обозначение на схеме & Значение/ Диапазон & Единица измерения & Шаг \\ 
\hline
\endfirsthead
\caption*{\hspace{3pt}\emph{Продолжение таблицы \ref{logzsign:tbl1}\hfill\vspace{-0.5\baselineskip}}} \\ % сделано по ГОСТ 2.105 п.6.8.7
\hline
\rowcolor{gray!20}
Параметр (Параметр на ИЧМ) & Условное обозначение на схеме & Значение/ Диапазон & Единица измерения & Шаг \\ 
%\hline
\endhead
%\hline
\endfoot
%\hline
\endlastfoot
%==+t1*PTRC1|TALMGASLGC
\multicolumn{5}{|c|}{ ЛО } \\ \hline 
\centering Ввод функции в работу (Ввод\_функции)  & \centering SGF1 & \centering Не предусмотрено \\ Предусмотрено & \centering -- & \centering \arraybackslash -- \\
\hline
\centering Выдержка времени срабатывания на блокировку (Тбл)  & \centering T1 & \centering 0...30000 & \centering мс & \centering \arraybackslash 10 \\
\hline
\centering Блокировка от низкой изоляции (Блок\_от\_НИ)  & \centering SGF2 & \centering Не предусмотрено \\ Предусмотрено & \centering -- & \centering \arraybackslash -- \\
\hline
%===t1
\end{longtable}
\normalsize


%%%%%%%%%%%%%%%%%%%%%%%%%%%%%%%%%%%%%%%%%%%%%%%%%%%%%%%%%% ЛО ГЗоткл %%%%%%%%%%%%%%%%%%%%%%%%%%%%%%%%%%%%%%%%%%%%%%%%%%%%%%%%%%%%%%%%%%
\par

В таблице \ref{logzotkl:tbl1} приведены параметры, необходимые для настройки ЛО ГЗоткл.

\footnotesize
\begin{longtable}{|>{\centering\arraybackslash}m{5.3cm}|>{\centering\arraybackslash}m{3.3cm}|>{\centering\arraybackslash}m{4.2cm}|>{\centering\arraybackslash}m{1.8cm}|>{\centering\arraybackslash}m{1cm}|}
\caption{Параметры для настройки логики отключения при срабатывании отключающего контакта газового реле трансформатора\hfill\vspace{-0.5\baselineskip}}\label{logzotkl:tbl1}\\ 
\hline
\rowcolor{gray!20}
Параметр (Параметр на ИЧМ) & Условное обозначение на схеме & Значение/ Диапазон & Единица измерения & Шаг \\ 
\hline
\endfirsthead
\caption*{\hspace{3pt}\emph{Продолжение таблицы \ref{logzotkl:tbl1}\hfill\vspace{-0.5\baselineskip}}} \\ % сделано по ГОСТ 2.105 п.6.8.7
\hline
\rowcolor{gray!20}
Параметр (Параметр на ИЧМ) & Условное обозначение на схеме & Значение/ Диапазон & Единица измерения & Шаг \\ 
%\hline
\endhead
%\hline
\endfoot
%\hline
\endlastfoot
%==+t1*PTRC1|TTRGASLGC
\multicolumn{5}{|c|}{ БУ } \\ \hline 
\centering Выдержка времени срабатывания на блокировку (Тбл)  & \centering T1 & \centering 0...30000 & \centering мс & \centering \arraybackslash 10 \\
\hline
\multicolumn{5}{|c|}{ ЛО } \\ \hline 
\centering Ввод функции в работу (Ввод\_функции)  & \centering SGF1 & \centering Не предусмотрено \\ Предусмотрено & \centering -- & \centering \arraybackslash -- \\
\hline
\centering Блокировка от низкой изоляции (Блок\_от\_НИ)  & \centering SGF2 & \centering Не предусмотрено \\ Предусмотрено & \centering -- & \centering \arraybackslash -- \\
\hline
%===t1
\end{longtable}
\normalsize


%%%%%%%%%%%%%%%%%%%%%%%%%%%%%%%%%%%%%%%%%%%%%%%%%%%%%%%%%% ЛО ТЗ %%%%%%%%%%%%%%%%%%%%%%%%%%%%%%%%%%%%%%%%%%%%%%%%%%%%%%%%%%%%%%%%%%
\par

В таблице \ref{lotz:tbl1} приведены параметры, необходимые для настройки ЛО ТЗ.
\footnotesize
\begin{longtable}{|>{\centering\arraybackslash}m{5.3cm}|>{\centering\arraybackslash}m{3.3cm}|>{\centering\arraybackslash}m{4.2cm}|>{\centering\arraybackslash}m{1.8cm}|>{\centering\arraybackslash}m{1cm}|}
\caption{Общие параметры для настройки ЛО ТЗ\hfill\vspace{-0.5\baselineskip}}\label{lotz:tbl1}\\ 
\hline
\rowcolor{gray!20}
Параметр (Параметр на ИЧМ) & Условное обозначение на схеме & Значение/ Диапазон & Единица измерения & Шаг \\ 
\hline
\endfirsthead
%\caption{\emph{Продолжение\hfill\vspace{-0.5\baselineskip}}} \\ % Отступ обозначения таблицы от таблицы и выравнивание надписи слева
\caption*{\hspace{3pt}\emph{Продолжение таблицы \ref{lotz:tbl1}\hfill\vspace{-0.5\baselineskip}}} \\ % сделано по ГОСТ 2.105 п.6.8.7
\hline
\rowcolor{gray!20}
Параметр (Параметр на ИЧМ) & Условное обозначение на схеме & Значение/ Диапазон & Единица измерения & Шаг \\ 
%\hline
\endhead
%\hline
\endfoot
%\hline
\endlastfoot
%==+t1*LLN0|APTTECHLGC
\multicolumn{5}{|c|}{ БУ } \\ \hline 
\centering Выдержка времени срабатывания на блокировку (Тбл)  & \centering T1 & \centering 0...30000 & \centering мс & \centering \arraybackslash 10 \\
\hline
\multicolumn{5}{|c|}{ ЛО ДТм } \\ \hline 
\centering Ввод функции в работу (Ввод\_функции)  & \centering SGF1 & \centering Не предусмотрено \\ Предусмотрено & \centering -- & \centering \arraybackslash -- \\
\hline
\centering Блокировка от низкой изоляции (Блок\_от\_НИ)  & \centering SGF2 & \centering Не предусмотрено \\ Предусмотрено & \centering -- & \centering \arraybackslash -- \\
\hline
\multicolumn{5}{|c|}{ ЛО ДТо } \\ \hline 
\centering Ввод функции в работу (Ввод\_функции)  & \centering SGF1 & \centering Не предусмотрено \\ Предусмотрено & \centering -- & \centering \arraybackslash -- \\
\hline
\centering Блокировка от низкой изоляции (Блок\_от\_НИ)  & \centering SGF2 & \centering Не предусмотрено \\ Предусмотрено & \centering -- & \centering \arraybackslash -- \\
\hline
\multicolumn{5}{|c|}{ ЛО РД } \\ \hline 
\centering Ввод функции в работу (Ввод\_функции)  & \centering SGF1 & \centering Не предусмотрено \\ Предусмотрено & \centering -- & \centering \arraybackslash -- \\
\hline
\centering Блокировка от низкой изоляции (Блок\_от\_НИ)  & \centering SGF2 & \centering Не предусмотрено \\ Предусмотрено & \centering -- & \centering \arraybackslash -- \\
\hline
%===t1
\end{longtable}
\normalsize


%%%%%%%%%%%%%%%%%%%%%%%%%%%%%%%%%%%%%%%%%%%%%%%%%%%%%%%%%% УРОВ %%%%%%%%%%%%%%%%%%%%%%%%%%%%%%%%%%%%%%%%%%%%%%%%%%%%%%%%%%%%%%%%%%
\par
В таблице \ref{urov35:tbl1} приведены параметры, необходимые для настройки УРОВ.

\footnotesize
\begin{longtable}{|>{\centering\arraybackslash}m{5.3cm}|>{\centering\arraybackslash}m{3.3cm}|>{\centering\arraybackslash}m{4.2cm}|>{\centering\arraybackslash}m{1.8cm}|>{\centering\arraybackslash}m{1cm}|}
\caption{Параметры для настройки УРОВ\hfill\vspace{-0.5\baselineskip}}\label{urov35:tbl1}\\ 
\hline
\rowcolor{gray!20}
Параметр (Параметр на ИЧМ) & Условное обозначение на схеме & Значение/ Диапазон & Единица измерения & Шаг \\ 
\hline
\endfirsthead
\caption*{\hspace{3pt}\emph{Продолжение таблицы \ref{urov35:tbl1}\hfill\vspace{-0.5\baselineskip}}} \\ % сделано по ГОСТ 2.105 п.6.8.7
\hline
\rowcolor{gray!20}
Параметр (Параметр на ИЧМ) & Условное обозначение на схеме & Значение/ Диапазон & Единица измерения & Шаг \\ 
%\hline
\endhead
%\hline
\endfoot
%\hline
\endlastfoot
%==+t1*GENRBRF1|TPBRF
\centering Ввод функции в работу (Ввод\_функции)  & \centering SGF1 & \centering Не предусмотрено \\ Предусмотрено & \centering -- & \centering \arraybackslash -- \\
\hline
\centering Ток срабатывания (Iср)  & \centering I> & \centering 0,05...0,50 & \centering о.е. & \centering \arraybackslash 0,01 \\
\hline
\centering Ускорение при блокировке отключения В (Уск\_при\_блк\_откл\_В)  & \centering SGF2 & \centering Не предусмотрено \\ Предусмотрено & \centering -- & \centering \arraybackslash -- \\
\hline
\centering УРОВ с подхватом по току (Подхват\_по\_току)  & \centering SGF3 & \centering Не предусмотрено \\ Предусмотрено & \centering -- & \centering \arraybackslash -- \\
\hline
\centering Контроль по току при действии "на себя" (Контр\_тока\_на\_себя)  & \centering SGF6 & \centering Не предусмотрено \\ Предусмотрено по внутр. ПО \\ Предусмотрено по внеш. ПО & \centering -- & \centering \arraybackslash -- \\
\hline
\centering Контроль ЭМО (Контроль\_ЭМО)  & \centering SGF4 & \centering Не предусмотрено \\ Предусмотрено по ЭМО1 \\ Предусмотрено по ЭМО1 и ЭМО2 & \centering -- & \centering \arraybackslash -- \\
\hline
\centering Действие внешнего УРОВ на вышестоящий выключатель (Действ\_на\_выш\_выкл)  & \centering SGF5 & \centering Не предусмотрено \\ Предусмотрено & \centering -- & \centering \arraybackslash -- \\
\hline
\centering Выдержка времени срабатывания (Tср)  & \centering T1 & \centering 0...1000 & \centering мс & \centering \arraybackslash 10 \\
\hline
%===t1
\end{longtable}
\normalsize


%%%%%%%%%%%%%%%%%%%%%%%%%%%%%%%%%%%%%%%%%%%%%%%%%%%%%%%%%% ЗМН1 %%%%%%%%%%%%%%%%%%%%%%%%%%%%%%%%%%%%%%%%%%%%%%%%%%%%%%%%%%%%%%%%%%
\par
В таблице \ref{zmn:tbl1} приведены параметры, необходимые для настройки ЗМН 1.

\footnotesize
\begin{longtable}{|>{\centering\arraybackslash}m{5.3cm}|>{\centering\arraybackslash}m{3.3cm}|>{\centering\arraybackslash}m{4.2cm}|>{\centering\arraybackslash}m{1.8cm}|>{\centering\arraybackslash}m{1cm}|}
\caption{Параметры для настройки ЗМН 1\hfill\vspace{-0.5\baselineskip}}\label{zmn:tbl1}\\ 
\hline
\rowcolor{gray!20}
Параметр (Параметр на ИЧМ) & Условное обозначение на схеме & Значение/ Диапазон & Единица измерения & Шаг \\ 
\hline
\endfirsthead
\caption*{\hspace{3pt}\emph{Продолжение таблицы \ref{zmn:tbl1}\hfill\vspace{-0.5\baselineskip}}} \\ % сделано по ГОСТ 2.105 п.6.8.7
\hline
\rowcolor{gray!20}
Параметр (Параметр на ИЧМ) & Условное обозначение на схеме & Значение/ Диапазон & Единица измерения & Шаг \\ 
%\hline
\endhead
%\hline
\endfoot
%\hline
\endlastfoot
%==+t1*PTUV1|LVTUV
\centering Ввод функции в работу (Ввод\_функции)  & \centering SGF1 & \centering Не предусмотрено \\ Предусмотрено & \centering -- & \centering \arraybackslash -- \\
\hline
\centering Напряжение срабатывания (Uср)  & \centering U< & \centering 2,0...100,0 & \centering В & \centering \arraybackslash 0,1 \\
\hline
\centering Напряжение срабатывания ОП (U2ср)  & \centering U2> & \centering 2,0...50,0 & \centering В & \centering \arraybackslash 0,1 \\
\hline
\centering Контроль напряжения ОП (Контроль\_U2)  & \centering SGF3 & \centering Не предусмотрено \\ Предусмотрено & \centering -- & \centering \arraybackslash -- \\
\hline
\centering Выбор режима работы (Реж\_работы)  & \centering SGF2 & \centering По 1 линейному \\ По 2 линейным \\ По 3 линейным & \centering -- & \centering \arraybackslash -- \\
\hline
\centering Выдержка времени срабатывания (Tср)  & \centering T1 & \centering 0...100000 & \centering мс & \centering \arraybackslash 10 \\
\hline
%===t1
\end{longtable}
\normalsize

%%%%%%%%%%%%%%%%%%%%%%%%%%%%%%%%%%%%%%%%%%%%%%%%%%%%%%%%%% ЗМН2 %%%%%%%%%%%%%%%%%%%%%%%%%%%%%%%%%%%%%%%%%%%%%%%%%%%%%%%%%%%%%%%%%%
\par
В таблице \ref{zmn1:tbl1} приведены параметры, необходимые для настройки ЗМН 2.

\footnotesize
\begin{longtable}{|>{\centering\arraybackslash}m{5.3cm}|>{\centering\arraybackslash}m{3.3cm}|>{\centering\arraybackslash}m{4.2cm}|>{\centering\arraybackslash}m{1.8cm}|>{\centering\arraybackslash}m{1cm}|}
\caption{Параметры для настройки ЗМН 2\hfill\vspace{-0.5\baselineskip}}\label{zmn1:tbl1}\\ 
\hline
\rowcolor{gray!20}
Параметр (Параметр на ИЧМ) & Условное обозначение на схеме & Значение/ Диапазон & Единица измерения & Шаг \\ 
\hline
\endfirsthead
\caption*{\hspace{3pt}\emph{Продолжение таблицы \ref{zmn1:tbl1}\hfill\vspace{-0.5\baselineskip}}} \\ % сделано по ГОСТ 2.105 п.6.8.7
\hline
\rowcolor{gray!20}
Параметр (Параметр на ИЧМ) & Условное обозначение на схеме & Значение/ Диапазон & Единица измерения & Шаг \\ 
%\hline
\endhead
%\hline
\endfoot
%\hline
\endlastfoot
%==+t1*PTUV2|LVTUV
\centering Ввод функции в работу (Ввод\_функции)  & \centering SGF1 & \centering Не предусмотрено \\ Предусмотрено & \centering -- & \centering \arraybackslash -- \\
\hline
\centering Напряжение срабатывания (Uср)  & \centering U< & \centering 2,0...100,0 & \centering В & \centering \arraybackslash 0,1 \\
\hline
\centering Напряжение срабатывания ОП (U2ср)  & \centering U2> & \centering 2,0...50,0 & \centering В & \centering \arraybackslash 0,1 \\
\hline
\centering Контроль напряжения ОП (Контроль\_U2)  & \centering SGF3 & \centering Не предусмотрено \\ Предусмотрено & \centering -- & \centering \arraybackslash -- \\
\hline
\centering Выбор режима работы (Реж\_работы)  & \centering SGF2 & \centering По 1 линейному \\ По 2 линейным \\ По 3 линейным & \centering -- & \centering \arraybackslash -- \\
\hline
\centering Выдержка времени срабатывания (Tср)  & \centering T1 & \centering 0...100000 & \centering мс & \centering \arraybackslash 10 \\
\hline
%===t1
\end{longtable}
\normalsize


%%%%%%%%%%%%%%%%%%%%%%%%%%%%%%%%%%%%%%%%%%%%%%%%%%%%%%%%%% ЗПН 1 %%%%%%%%%%%%%%%%%%%%%%%%%%%%%%%%%%%%%%%%%%%%%%%%%%%%%%%%%%%%%%%%%%
\par
В таблице \ref{zpn:tbl1} приведены параметры, необходимые для настройки ЗПН 1.

\footnotesize
\begin{longtable}{|>{\centering\arraybackslash}m{5.3cm}|>{\centering\arraybackslash}m{3.3cm}|>{\centering\arraybackslash}m{4.2cm}|>{\centering\arraybackslash}m{1.8cm}|>{\centering\arraybackslash}m{1cm}|}
\caption{Параметры для настройки ЗПН 1\hfill\vspace{-0.5\baselineskip}}\label{zpn:tbl1}\\ 
\hline
\rowcolor{gray!20}
Параметр (Параметр на ИЧМ) & Условное обозначение на схеме & Значение/ Диапазон & Единица измерения & Шаг \\ 
\hline
\endfirsthead
\caption*{\hspace{3pt}\emph{Продолжение таблицы \ref{zpn:tbl1}\hfill\vspace{-0.5\baselineskip}}} \\ % сделано по ГОСТ 2.105 п.6.8.7
\hline
\rowcolor{gray!20}
Параметр (Параметр на ИЧМ) & Условное обозначение на схеме & Значение/ Диапазон & Единица измерения & Шаг \\ 
%\hline
\endhead
%\hline
\endfoot
%\hline
\endlastfoot
%==+t1*PTOV1|LVTOV
\centering Ввод функции в работу (Ввод\_функции)  & \centering SGF1 & \centering Не предусмотрено \\ Предусмотрено & \centering -- & \centering \arraybackslash -- \\
\hline
\centering Напряжение срабатывания (Uср)  & \centering U> & \centering 50,0...250,0 & \centering В & \centering \arraybackslash 0,1 \\
\hline
\centering Напряжение срабатывания ОП (U2ср)  & \centering U2> & \centering 2,0...50,0 & \centering В & \centering \arraybackslash 0,1 \\
\hline
\centering Контроль напряжения ОП (Контроль\_U2)  & \centering SGF3 & \centering Не предусмотрено \\ Предусмотрено & \centering -- & \centering \arraybackslash -- \\
\hline
\centering Выбор режима работы (Реж\_работы)  & \centering SGF2 & \centering По 1 линейному \\ По 2 линейным \\ По 3 линейным & \centering -- & \centering \arraybackslash -- \\
\hline
\centering Выдержка времени срабатывания (Tср)  & \centering T1 & \centering 0...100000 & \centering мс & \centering \arraybackslash 10 \\
\hline
%===t1
\end{longtable}
\normalsize


%%%%%%%%%%%%%%%%%%%%%%%%%%%%%%%%%%%%%%%%%%%%%%%%%%%%%%%%%% ЗПН 2 %%%%%%%%%%%%%%%%%%%%%%%%%%%%%%%%%%%%%%%%%%%%%%%%%%%%%%%%%%%%%%%%%%
\par
В таблице \ref{zpn1:tbl1} приведены параметры, необходимые для настройки ЗПН 2.

\footnotesize
\begin{longtable}{|>{\centering\arraybackslash}m{5.3cm}|>{\centering\arraybackslash}m{3.3cm}|>{\centering\arraybackslash}m{4.2cm}|>{\centering\arraybackslash}m{1.8cm}|>{\centering\arraybackslash}m{1cm}|}
\caption{Параметры для настройки ЗПН 2\hfill\vspace{-0.5\baselineskip}}\label{zpn1:tbl1}\\ 
\hline
\rowcolor{gray!20}
Параметр (Параметр на ИЧМ) & Условное обозначение на схеме & Значение/ Диапазон & Единица измерения & Шаг \\ 
\hline
\endfirsthead
\caption*{\hspace{3pt}\emph{Продолжение таблицы \ref{zpn1:tbl1}\hfill\vspace{-0.5\baselineskip}}} \\ % сделано по ГОСТ 2.105 п.6.8.7
\hline
\rowcolor{gray!20}
Параметр (Параметр на ИЧМ) & Условное обозначение на схеме & Значение/ Диапазон & Единица измерения & Шаг \\ 
%\hline
\endhead
%\hline
\endfoot
%\hline
\endlastfoot
%==+t1*PTOV2|LVTOV
\centering Ввод функции в работу (Ввод\_функции)  & \centering SGF1 & \centering Не предусмотрено \\ Предусмотрено & \centering -- & \centering \arraybackslash -- \\
\hline
\centering Напряжение срабатывания (Uср)  & \centering U> & \centering 50,0...250,0 & \centering В & \centering \arraybackslash 0,1 \\
\hline
\centering Напряжение срабатывания ОП (U2ср)  & \centering U2> & \centering 2,0...50,0 & \centering В & \centering \arraybackslash 0,1 \\
\hline
\centering Контроль напряжения ОП (Контроль\_U2)  & \centering SGF3 & \centering Не предусмотрено \\ Предусмотрено & \centering -- & \centering \arraybackslash -- \\
\hline
\centering Выбор режима работы (Реж\_работы)  & \centering SGF2 & \centering По 1 линейному \\ По 2 линейным \\ По 3 линейным & \centering -- & \centering \arraybackslash -- \\
\hline
\centering Выдержка времени срабатывания (Tср)  & \centering T1 & \centering 0...100000 & \centering мс & \centering \arraybackslash 10 \\
\hline
%===t1
\end{longtable}
\normalsize



%%%%%%%%%%%%%%%%%%%%%%%%%%%%%%%%%%%%%%%%%%%%%%%%%%%%%%%%%% БННш %%%%%%%%%%%%%%%%%%%%%%%%%%%%%%%%%%%%%%%%%%%%%%%%%%%%%%%%%%%%%%%%%%
\par
В таблице \ref{bnnsh:tbl1} приведены параметры, необходимые для настройки БННш.

\footnotesize
\begin{longtable}{|>{\centering\arraybackslash}m{5.3cm}|>{\centering\arraybackslash}m{3.3cm}|>{\centering\arraybackslash}m{4.2cm}|>{\centering\arraybackslash}m{1.8cm}|>{\centering\arraybackslash}m{1cm}|}
\caption{Параметры для настройки БННш\hfill\vspace{-0.5\baselineskip}}\label{bnnsh:tbl1}\\ 
\hline
\rowcolor{gray!20}
Параметр (Параметр на ИЧМ) & Условное обозначение на схеме & Значение/ Диапазон & Единица измерения & Шаг \\ 
\hline
\endfirsthead
\caption*{\hspace{3pt}\emph{Продолжение таблицы \ref{bnnsh:tbl1}\hfill\vspace{-0.5\baselineskip}}} \\ % сделано по ГОСТ 2.105 п.6.8.7
\hline
\rowcolor{gray!20}
Параметр (Параметр на ИЧМ) & Условное обозначение на схеме & Значение/ Диапазон & Единица измерения & Шаг \\ 
%\hline
\endhead
%\hline
\endfoot
%\hline
\endlastfoot
%==+t1*RVTR1|LVVTR1
\centering Ввод функции в работу (Ввод\_функции)  & \centering SGF1 & \centering Не предусмотрено \\ Предусмотрено & \centering -- & \centering \arraybackslash -- \\
\hline
\centering Уровень небаланса по напряжению (dU0ср)  & \centering ∑Uф--3U0> & \centering 1,0...150,0 & \centering В & \centering \arraybackslash 0,1 \\
\hline
\centering Ток срабатывания ОП (I2ср)  & \centering I2< & \centering 0,04...4,00 & \centering о.е. & \centering \arraybackslash 0,01 \\
\hline
\centering Напряжение срабатывания ОП (U2ср)  & \centering U2> & \centering 2,0...60,0 & \centering В & \centering \arraybackslash 0,1 \\
\hline
\centering Напряжение срабатывания (Uср)  & \centering U< & \centering 2,0...100,0 & \centering В & \centering \arraybackslash 0,1 \\
\hline
\centering Ток срабатывания (Iср)  & \centering I> & \centering 0,04...4,00 & \centering о.е. & \centering \arraybackslash 0,01 \\
\hline
\centering Напряжение срабатывания НП третьей гармоники (3U0ср\_3гарм)  & \centering 3U03f< & \centering 0,05...10,00 & \centering В & \centering \arraybackslash 0,01 \\
\hline
\centering Ввод в работу контроля небаланса (Контроль\_dU0)  & \centering SGF2 & \centering Не предусмотрено \\ Предусмотрено & \centering -- & \centering \arraybackslash -- \\
\hline
\centering Ввод в работу контроля напряжения ОП (Контроль\_U2)  & \centering SGF3 & \centering Не предусмотрено \\ Предусмотрено & \centering -- & \centering \arraybackslash -- \\
\hline
\centering Ввод в работу контроля минимального напряжения (Контроль\_Uмин)  & \centering SGF4 & \centering Не предусмотрено \\ Предусмотрено & \centering -- & \centering \arraybackslash -- \\
\hline
\centering Ввод в работу контроля 3 гармоники напряжения НП (Контроль\_3гарм\_3U0)  & \centering SGF5 & \centering Не предусмотрено \\ Предусмотрено & \centering -- & \centering \arraybackslash -- \\
\hline
\centering Выдержка времени срабатывания сигнализации (Tсигн)  & \centering T1 & \centering 0...100000 & \centering мс & \centering \arraybackslash 10 \\
\hline
%===t1
\end{longtable}
\normalsize


%%%%%%%%%%%%%%%%%%%%%%%%%%%%%%%%%%%%%%%%%%%%%%%%%%%%%%%%%% БННп %%%%%%%%%%%%%%%%%%%%%%%%%%%%%%%%%%%%%%%%%%%%%%%%%%%%%%%%%%%%%%%%%%
\par
В таблице \ref{bnnp:tbl1} приведены параметры, необходимые для настройки БННп.

\footnotesize
\begin{longtable}{|>{\centering\arraybackslash}m{5.3cm}|>{\centering\arraybackslash}m{3.3cm}|>{\centering\arraybackslash}m{4.2cm}|>{\centering\arraybackslash}m{1.8cm}|>{\centering\arraybackslash}m{1cm}|}
\caption{Параметры для настройки БННп\hfill\vspace{-0.5\baselineskip}}\label{bnnp:tbl1}\\ 
\hline
\rowcolor{gray!20}
Параметр (Параметр на ИЧМ) & Условное обозначение на схеме & Значение/ Диапазон & Единица измерения & Шаг \\ 
\hline
\endfirsthead
\caption*{\hspace{3pt}\emph{Продолжение таблицы \ref{bnnp:tbl1}\hfill\vspace{-0.5\baselineskip}}} \\ % сделано по ГОСТ 2.105 п.6.8.7
\hline
\rowcolor{gray!20}
Параметр (Параметр на ИЧМ) & Условное обозначение на схеме & Значение/ Диапазон & Единица измерения & Шаг \\ 
%\hline
\endhead
%\hline
\endfoot
%\hline
\endlastfoot
%==+t1*RVTR1|LVVTR2
\centering Ввод функции в работу (Ввод\_функции)  & \centering SGF1 & \centering Не предусмотрено \\ Предусмотрено & \centering -- & \centering \arraybackslash -- \\
\hline
\centering Ток срабатывания ОП (I2ср)  & \centering I2< & \centering 0,04...4,00 & \centering о.е. & \centering \arraybackslash 0,01 \\
\hline
\centering Напряжение срабатывания ОП (U2ср)  & \centering U2> & \centering 2,0...60,0 & \centering В & \centering \arraybackslash 0,1 \\
\hline
\centering Напряжение срабатывания (Uср)  & \centering U< & \centering 2,0...100,0 & \centering В & \centering \arraybackslash 0,1 \\
\hline
\centering Ток срабатывания (Iср)  & \centering I> & \centering 0,04...4,00 & \centering о.е. & \centering \arraybackslash 0,01 \\
\hline
\centering Ввод в работу контроля напряжения ОП (Контроль\_U2)  & \centering SGF2 & \centering Не предусмотрено \\ Предусмотрено & \centering -- & \centering \arraybackslash -- \\
\hline
\centering Ввод в работу контроля минимального напряжения (Контроль\_Uмин)  & \centering SGF3 & \centering Не предусмотрено \\ Предусмотрено & \centering -- & \centering \arraybackslash -- \\
\hline
\centering Выдержка времени срабатывания сигнализации (Tсигн)  & \centering T1 & \centering 0...100000 & \centering мс & \centering \arraybackslash 10 \\
\hline
%===t1
\end{longtable}
\normalsize



%%%%%%%%%%%%%%%%%%%%%%%%%%%%%%%%%%%%%%%%%%%%%%%%%%%%%%%%%% ЗСЧ 1 %%%%%%%%%%%%%%%%%%%%%%%%%%%%%%%%%%%%%%%%%%%%%%%%%%%%%%%%%%%%%%%%%%
\par
В таблице \ref{zsch:tbl1} приведены параметры, необходимые для настройки ЗСЧ 1.

\footnotesize
\begin{longtable}{|>{\centering\arraybackslash}m{5.3cm}|>{\centering\arraybackslash}m{3.3cm}|>{\centering\arraybackslash}m{4.2cm}|>{\centering\arraybackslash}m{1.8cm}|>{\centering\arraybackslash}m{1cm}|}
\caption{Параметры для настройки ЗСЧ 1\hfill\vspace{-0.5\baselineskip}}\label{zsch:tbl1}\\ 
\hline
\rowcolor{gray!20}
Параметр (Параметр на ИЧМ) & Условное обозначение на схеме & Значение/ Диапазон & Единица измерения & Шаг \\ 
\hline
\endfirsthead
\caption*{\hspace{3pt}\emph{Продолжение таблицы \ref{zsch:tbl1}\hfill\vspace{-0.5\baselineskip}}} \\ % сделано по ГОСТ 2.105 п.6.8.7
\hline
\rowcolor{gray!20}
Параметр (Параметр на ИЧМ) & Условное обозначение на схеме & Значение/ Диапазон & Единица измерения & Шаг \\ 
%\hline
\endhead
%\hline
\endfoot
%\hline
\endlastfoot
%==+t1*PTUF1|LVTUF
\centering Ввод функции в работу (Ввод\_функции)  & \centering SGF1 & \centering Не предусмотрено \\ Предусмотрено & \centering -- & \centering \arraybackslash -- \\
\hline
\centering Частота срабатывания (Fср)  & \centering F< & \centering 45,0...51,0 & \centering Гц & \centering \arraybackslash 0,1 \\
\hline
\centering Разность частот срабатывания и возврата (dFвоз)  & \centering -- & \centering 0,04...1,00 & \centering Гц & \centering \arraybackslash 0,01 \\
\hline
\centering Напряжение срабатывания (Uср)  & \centering U< & \centering 2,0...100,0 & \centering В & \centering \arraybackslash 0,1 \\
\hline
\centering Величина срабатывания по скорости изменения частоты (dFdtср)  & \centering dF</dt & \centering 0,1...15,0 & \centering Гц/с & \centering \arraybackslash 0,1 \\
\hline
\centering Коэффициент возврата по скорости изменения частоты (Kвоз)  & \centering -- & \centering 0,20...0,99 & \centering -- & \centering \arraybackslash 0,01 \\
\hline
\centering Выдержка времени на возврат (Tвоз)  & \centering T2 & \centering 0...10000 & \centering мс & \centering \arraybackslash 10 \\
\hline
\centering Блокировка по dF/dt (Блок\_по\_dF/dt)  & \centering SGF2 & \centering Не предусмотрено \\ Предусмотрено & \centering -- & \centering \arraybackslash -- \\
\hline
\centering Выдержка времени срабатывания (Tср)  & \centering T1 & \centering 0...100000 & \centering мс & \centering \arraybackslash 10 \\
\hline
%===t1
\end{longtable}
\normalsize


%%%%%%%%%%%%%%%%%%%%%%%%%%%%%%%%%%%%%%%%%%%%%%%%%%%%%%%%%% ЗСЧ 2 %%%%%%%%%%%%%%%%%%%%%%%%%%%%%%%%%%%%%%%%%%%%%%%%%%%%%%%%%%%%%%%%%%
\par
В таблице \ref{zsch1:tbl1} приведены параметры, необходимые для настройки ЗСЧ 2.

\footnotesize
\begin{longtable}{|>{\centering\arraybackslash}m{5.3cm}|>{\centering\arraybackslash}m{3.3cm}|>{\centering\arraybackslash}m{4.2cm}|>{\centering\arraybackslash}m{1.8cm}|>{\centering\arraybackslash}m{1cm}|}
\caption{Параметры для настройки ЗСЧ 2\hfill\vspace{-0.5\baselineskip}}\label{zsch1:tbl1}\\ 
\hline
\rowcolor{gray!20}
Параметр (Параметр на ИЧМ) & Условное обозначение на схеме & Значение/ Диапазон & Единица измерения & Шаг \\ 
\hline
\endfirsthead
\caption*{\hspace{3pt}\emph{Продолжение таблицы \ref{zsch1:tbl1}\hfill\vspace{-0.5\baselineskip}}} \\ % сделано по ГОСТ 2.105 п.6.8.7
\hline
\rowcolor{gray!20}
Параметр (Параметр на ИЧМ) & Условное обозначение на схеме & Значение/ Диапазон & Единица измерения & Шаг \\ 
%\hline
\endhead
%\hline
\endfoot
%\hline
\endlastfoot
%==+t1*PTUF2|LVTUF
\centering Ввод функции в работу (Ввод\_функции)  & \centering SGF1 & \centering Не предусмотрено \\ Предусмотрено & \centering -- & \centering \arraybackslash -- \\
\hline
\centering Частота срабатывания (Fср)  & \centering F< & \centering 45,0...51,0 & \centering Гц & \centering \arraybackslash 0,1 \\
\hline
\centering Разность частот срабатывания и возврата (dFвоз)  & \centering -- & \centering 0,04...1,00 & \centering Гц & \centering \arraybackslash 0,01 \\
\hline
\centering Напряжение срабатывания (Uср)  & \centering U< & \centering 2,0...100,0 & \centering В & \centering \arraybackslash 0,1 \\
\hline
\centering Величина срабатывания по скорости изменения частоты (dFdtср)  & \centering dF</dt & \centering 0,1...15,0 & \centering Гц/с & \centering \arraybackslash 0,1 \\
\hline
\centering Коэффициент возврата по скорости изменения частоты (Kвоз)  & \centering -- & \centering 0,20...0,99 & \centering -- & \centering \arraybackslash 0,01 \\
\hline
\centering Выдержка времени на возврат (Tвоз)  & \centering T2 & \centering 0...10000 & \centering мс & \centering \arraybackslash 10 \\
\hline
\centering Блокировка по dF/dt (Блок\_по\_dF/dt)  & \centering SGF2 & \centering Не предусмотрено \\ Предусмотрено & \centering -- & \centering \arraybackslash -- \\
\hline
\centering Выдержка времени срабатывания (Tср)  & \centering T1 & \centering 0...100000 & \centering мс & \centering \arraybackslash 10 \\
\hline
%===t1
\end{longtable}
\normalsize


%%%%%%%%%%%%%%%%%%%%%%%%%%%%%%%%%%%%%%%%%%%%%%%%%%%%%%%%%% ЗПЧ 1 %%%%%%%%%%%%%%%%%%%%%%%%%%%%%%%%%%%%%%%%%%%%%%%%%%%%%%%%%%%%%%%%%%
\par
В таблице \ref{zpch:tbl1} приведены параметры, необходимые для настройки ЗПЧ 1.

\footnotesize
\begin{longtable}{|>{\centering\arraybackslash}m{5.3cm}|>{\centering\arraybackslash}m{3.3cm}|>{\centering\arraybackslash}m{4.2cm}|>{\centering\arraybackslash}m{1.8cm}|>{\centering\arraybackslash}m{1cm}|}
\caption{Параметры для настройки ЗПЧ 1\hfill\vspace{-0.5\baselineskip}}\label{zpch:tbl1}\\ 
\hline
\rowcolor{gray!20}
Параметр (Параметр на ИЧМ) & Условное обозначение на схеме & Значение/ Диапазон & Единица измерения & Шаг \\ 
\hline
\endfirsthead
\caption*{\hspace{3pt}\emph{Продолжение таблицы \ref{zpch:tbl1}\hfill\vspace{-0.5\baselineskip}}} \\ % сделано по ГОСТ 2.105 п.6.8.7
\hline
\rowcolor{gray!20}
Параметр (Параметр на ИЧМ) & Условное обозначение на схеме & Значение/ Диапазон & Единица измерения & Шаг \\ 
%\hline
\endhead
%\hline
\endfoot
%\hline
\endlastfoot
%==+t1*PTOF1|LVTOF
\centering Ввод функции в работу (Ввод\_функции)  & \centering SGF1 & \centering Не предусмотрено \\ Предусмотрено & \centering -- & \centering \arraybackslash -- \\
\hline
\centering Минимальное рабочее напряжение (Uср)  & \centering U< & \centering 2,0...100,0 & \centering В & \centering \arraybackslash 0,1 \\
\hline
\centering Уровень частоты срабатывания (Fср)  & \centering F> & \centering 49,0...55,0 & \centering Гц & \centering \arraybackslash 0,1 \\
\hline
\centering Разность частот срабатывания и возврата (dFвоз)  & \centering -- & \centering 0,04...1,00 & \centering Гц & \centering \arraybackslash 0,01 \\
\hline
\centering Выдержка времени на возврат (Tвоз)  & \centering T2 & \centering 0...10000 & \centering мс & \centering \arraybackslash 10 \\
\hline
\centering Выдержка времени срабатывания (Tср)  & \centering T1 & \centering 0...100000 & \centering мс & \centering \arraybackslash 10 \\
\hline
%===t1
\end{longtable}
\normalsize


%%%%%%%%%%%%%%%%%%%%%%%%%%%%%%%%%%%%%%%%%%%%%%%%%%%%%%%%%% ЗПЧ 2 %%%%%%%%%%%%%%%%%%%%%%%%%%%%%%%%%%%%%%%%%%%%%%%%%%%%%%%%%%%%%%%%%%
\par
В таблице \ref{zpch1:tbl1} приведены параметры, необходимые для настройки ЗПЧ 2.

\footnotesize
\begin{longtable}{|>{\centering\arraybackslash}m{5.3cm}|>{\centering\arraybackslash}m{3.3cm}|>{\centering\arraybackslash}m{4.2cm}|>{\centering\arraybackslash}m{1.8cm}|>{\centering\arraybackslash}m{1cm}|}
\caption{Параметры для настройки ЗПЧ 2\hfill\vspace{-0.5\baselineskip}}\label{zpch1:tbl1}\\ 
\hline
\rowcolor{gray!20}
Параметр (Параметр на ИЧМ) & Условное обозначение на схеме & Значение/ Диапазон & Единица измерения & Шаг \\ 
\hline
\endfirsthead
\caption*{\hspace{3pt}\emph{Продолжение таблицы \ref{zpch1:tbl1}\hfill\vspace{-0.5\baselineskip}}} \\ % сделано по ГОСТ 2.105 п.6.8.7
\hline
\rowcolor{gray!20}
Параметр (Параметр на ИЧМ) & Условное обозначение на схеме & Значение/ Диапазон & Единица измерения & Шаг \\ 
%\hline
\endhead
%\hline
\endfoot
%\hline
\endlastfoot
%==+t1*PTOF2|LVTOF
\centering Ввод функции в работу (Ввод\_функции)  & \centering SGF1 & \centering Не предусмотрено \\ Предусмотрено & \centering -- & \centering \arraybackslash -- \\
\hline
\centering Минимальное рабочее напряжение (Uср)  & \centering U< & \centering 2,0...100,0 & \centering В & \centering \arraybackslash 0,1 \\
\hline
\centering Уровень частоты срабатывания (Fср)  & \centering F> & \centering 49,0...55,0 & \centering Гц & \centering \arraybackslash 0,1 \\
\hline
\centering Разность частот срабатывания и возврата (dFвоз)  & \centering -- & \centering 0,04...1,00 & \centering Гц & \centering \arraybackslash 0,01 \\
\hline
\centering Выдержка времени на возврат (Tвоз)  & \centering T2 & \centering 0...10000 & \centering мс & \centering \arraybackslash 10 \\
\hline
\centering Выдержка времени срабатывания (Tср)  & \centering T1 & \centering 0...100000 & \centering мс & \centering \arraybackslash 10 \\
\hline
%===t1
\end{longtable}
\normalsize


%%%%%%%%%%%%%%%%%%%%%%%%%%%%%%%%%%%%%%%%%%%%%%%%%%%%%%%%%% ЗИЧ 1 %%%%%%%%%%%%%%%%%%%%%%%%%%%%%%%%%%%%%%%%%%%%%%%%%%%%%%%%%%%%%%%%%%
\par
В таблице \ref{zich:tbl1} приведены параметры, необходимые для настройки ЗИЧ 1.

\footnotesize
\begin{longtable}{|>{\centering\arraybackslash}m{5.3cm}|>{\centering\arraybackslash}m{3.3cm}|>{\centering\arraybackslash}m{4.2cm}|>{\centering\arraybackslash}m{1.8cm}|>{\centering\arraybackslash}m{1cm}|}
\caption{Параметры для настройки ЗИЧ 1\hfill\vspace{-0.5\baselineskip}}\label{zich:tbl1}\\ 
\hline
\rowcolor{gray!20}
Параметр (Параметр на ИЧМ) & Условное обозначение на схеме & Значение/ Диапазон & Единица измерения & Шаг \\ 
\hline
\endfirsthead
\caption*{\hspace{3pt}\emph{Продолжение таблицы \ref{zich:tbl1}\hfill\vspace{-0.5\baselineskip}}} \\ % сделано по ГОСТ 2.105 п.6.8.7
\hline
\rowcolor{gray!20}
Параметр (Параметр на ИЧМ) & Условное обозначение на схеме & Значение/ Диапазон & Единица измерения & Шаг \\ 
%\hline
\endhead
%\hline
\endfoot
%\hline
\endlastfoot
%==+t1*BLKPFRC1|LVPFRC
\centering Ввод функции в работу (Ввод\_функции)  & \centering SGF1 & \centering Не предусмотрено \\ Предусмотрено & \centering -- & \centering \arraybackslash -- \\
\hline
\centering Величина срабатывания по скорости изменения частоты (dFdtСр)  & \centering dF/dt & \centering 0,1...15,0 & \centering Гц/с & \centering \arraybackslash 0,1 \\
\hline
\centering Коэффициент возврата по скорости изменения частоты (Kвоз)  & \centering -- & \centering 0,20...0,99 & \centering -- & \centering \arraybackslash 0,01 \\
\hline
\centering Напряжение срабатывания (Uср)  & \centering U< & \centering 2,0...100,0 & \centering В & \centering \arraybackslash 0,1 \\
\hline
\centering Режим работы ЗИЧ (Реж\_ЗИЧ)  & \centering SGF2 & \centering по dF>/dt \\ по dF</dt & \centering -- & \centering \arraybackslash -- \\
\hline
\centering Выдержка времени на возврат (Tвоз)  & \centering T2 & \centering 0...10000 & \centering мс & \centering \arraybackslash 10 \\
\hline
\centering Выдержка времени срабатывания (Tср)  & \centering T1 & \centering 0...100000 & \centering мс & \centering \arraybackslash 10 \\
\hline
%===t1
\end{longtable}
\normalsize


%%%%%%%%%%%%%%%%%%%%%%%%%%%%%%%%%%%%%%%%%%%%%%%%%%%%%%%%%% ЗИЧ 2 %%%%%%%%%%%%%%%%%%%%%%%%%%%%%%%%%%%%%%%%%%%%%%%%%%%%%%%%%%%%%%%%%%
\par
В таблице \ref{zich1:tbl1} приведены параметры, необходимые для настройки ЗИЧ 2.

\footnotesize
\begin{longtable}{|>{\centering\arraybackslash}m{5.3cm}|>{\centering\arraybackslash}m{3.3cm}|>{\centering\arraybackslash}m{4.2cm}|>{\centering\arraybackslash}m{1.8cm}|>{\centering\arraybackslash}m{1cm}|}
\caption{Параметры для настройки ЗИЧ 2\hfill\vspace{-0.5\baselineskip}}\label{zich1:tbl1}\\ 
\hline
\rowcolor{gray!20}
Параметр (Параметр на ИЧМ) & Условное обозначение на схеме & Значение/ Диапазон & Единица измерения & Шаг \\ 
\hline
\endfirsthead
\caption*{\hspace{3pt}\emph{Продолжение таблицы \ref{zich1:tbl1}\hfill\vspace{-0.5\baselineskip}}} \\ % сделано по ГОСТ 2.105 п.6.8.7
\hline
\rowcolor{gray!20}
Параметр (Параметр на ИЧМ) & Условное обозначение на схеме & Значение/ Диапазон & Единица измерения & Шаг \\ 
%\hline
\endhead
%\hline
\endfoot
%\hline
\endlastfoot
%==+t1*BLKPFRC2|LVPFRC
\centering Ввод функции в работу (Ввод\_функции)  & \centering SGF1 & \centering Не предусмотрено \\ Предусмотрено & \centering -- & \centering \arraybackslash -- \\
\hline
\centering Величина срабатывания по скорости изменения частоты (dFdtСр)  & \centering dF/dt & \centering 0,1...15,0 & \centering Гц/с & \centering \arraybackslash 0,1 \\
\hline
\centering Коэффициент возврата по скорости изменения частоты (Kвоз)  & \centering -- & \centering 0,20...0,99 & \centering -- & \centering \arraybackslash 0,01 \\
\hline
\centering Напряжение срабатывания (Uср)  & \centering U< & \centering 2,0...100,0 & \centering В & \centering \arraybackslash 0,1 \\
\hline
\centering Режим работы ЗИЧ (Реж\_ЗИЧ)  & \centering SGF2 & \centering по dF>/dt \\ по dF</dt & \centering -- & \centering \arraybackslash -- \\
\hline
\centering Выдержка времени на возврат (Tвоз)  & \centering T2 & \centering 0...10000 & \centering мс & \centering \arraybackslash 10 \\
\hline
\centering Выдержка времени срабатывания (Tср)  & \centering T1 & \centering 0...100000 & \centering мс & \centering \arraybackslash 10 \\
\hline
%===t1
\end{longtable}
\normalsize


%%%%%%%%%%%%%%%%%%%%%%%%%%%%%%%%%%%%%%%%%%%%%%%%%%%%%%%%%% АЧР %%%%%%%%%%%%%%%%%%%%%%%%%%%%%%%%%%%%%%%%%%%%%%%%%%%%%%%%%%%%%%%%%%
\par
В таблице \ref{achr:tbl1} приведены параметры, необходимые для настройки АЧР.

\footnotesize
\begin{longtable}{|>{\centering\arraybackslash}m{5.3cm}|>{\centering\arraybackslash}m{3.3cm}|>{\centering\arraybackslash}m{4.2cm}|>{\centering\arraybackslash}m{1.8cm}|>{\centering\arraybackslash}m{1cm}|}
\caption{Параметры для настройки АЧР\hfill\vspace{-0.5\baselineskip}}\label{achr:tbl1}\\ 
\hline
\rowcolor{gray!20}
Параметр (Параметр на ИЧМ) & Условное обозначение на схеме & Значение/ Диапазон & Единица измерения & Шаг \\ 
\hline
\endfirsthead
\caption*{\hspace{3pt}\emph{Продолжение таблицы \ref{achr:tbl1}\hfill\vspace{-0.5\baselineskip}}} \\ % сделано по ГОСТ 2.105 п.6.8.7
\hline
\rowcolor{gray!20}
Параметр (Параметр на ИЧМ) & Условное обозначение на схеме & Значение/ Диапазон & Единица измерения & Шаг \\ 
%\hline
\endhead
%\hline
\endfoot
%\hline
\endlastfoot
%==+t1*ECAPTUF1|ECATUF
\centering Ввод функции в работу (Ввод\_функции)  & \centering SGF1 & \centering Не предусмотрено \\ Предусмотрено & \centering -- & \centering \arraybackslash -- \\
\hline
\centering Частота срабатывания (Fср)  & \centering F< & \centering 45,0...51,0 & \centering Гц & \centering \arraybackslash 0,1 \\
\hline
\centering Разность частот срабатывания и возврата (dFвозв)  & \centering -- & \centering 0,04...1,00 & \centering Гц & \centering \arraybackslash 0,01 \\
\hline
\centering Напряжение срабатывания (Uср)  & \centering U< & \centering 2,0...100,0 & \centering В & \centering \arraybackslash 0,1 \\
\hline
\centering Величина срабатывания по скорости изменения частоты (dFdtср)  & \centering dF</dt & \centering 0,1...15,0 & \centering Гц/с & \centering \arraybackslash 0,1 \\
\hline
\centering Коэффициент возврата по скорости изменения частоты (Kвоз)  & \centering -- & \centering 0,20...0,99 & \centering -- & \centering \arraybackslash 0,01 \\
\hline
\centering Выдержка времени на возврат (Tвоз)  & \centering T2 & \centering 0...10000 & \centering мс & \centering \arraybackslash 10 \\
\hline
\centering Блокировка по dF/dt (Блок\_по\_dF/dt)  & \centering SGF3 & \centering Не предусмотрено \\ Предусмотрено & \centering -- & \centering \arraybackslash -- \\
\hline
\centering Режим работы (Реж\_работы)  & \centering SGF2 & \centering По внешнему сигналу \\ Измеренный & \centering -- & \centering \arraybackslash -- \\
\hline
\centering Выдержка времени срабатывания (Tср)  & \centering T1 & \centering 0...100000 & \centering мс & \centering \arraybackslash 10 \\
\hline
\centering Длительность импульса (Tимп)  & \centering T3 & \centering 50...60000 & \centering мс & \centering \arraybackslash 10 \\
\hline
%===t1
\end{longtable}
\normalsize


%%%%%%%%%%%%%%%%%%%%%%%%%%%%%%%%%%%%%%%%%%%%%%%%%%%%%%%%%% ЧАПВ %%%%%%%%%%%%%%%%%%%%%%%%%%%%%%%%%%%%%%%%%%%%%%%%%%%%%%%%%%%%%%%%%%
\par
В таблице \ref{chapv:tbl1} приведены параметры, необходимые для настройки ЧАПВ.

\footnotesize
\begin{longtable}{|>{\centering\arraybackslash}m{5.3cm}|>{\centering\arraybackslash}m{3.3cm}|>{\centering\arraybackslash}m{4.2cm}|>{\centering\arraybackslash}m{1.8cm}|>{\centering\arraybackslash}m{1cm}|}
\caption{Параметры для настройки ЧАПВ\hfill\vspace{-0.5\baselineskip}}\label{chapv:tbl1}\\ 
\hline
\rowcolor{gray!20}
Параметр (Параметр на ИЧМ) & Условное обозначение на схеме & Значение/ Диапазон & Единица измерения & Шаг \\ 
\hline
\endfirsthead
\caption*{\hspace{3pt}\emph{Продолжение таблицы \ref{chapv:tbl1}\hfill\vspace{-0.5\baselineskip}}} \\ % сделано по ГОСТ 2.105 п.6.8.7
\hline
\rowcolor{gray!20}
Параметр (Параметр на ИЧМ) & Условное обозначение на схеме & Значение/ Диапазон & Единица измерения & Шаг \\ 
%\hline
\endhead
%\hline
\endfoot
%\hline
\endlastfoot
%==+t1*ECAPTOF1|LVECATOF
\centering Ввод функции в работу (Ввод\_функции)  & \centering SGF1 & \centering Не предусмотрено \\ Предусмотрено & \centering -- & \centering \arraybackslash -- \\
\hline
\centering Блокировка ЧАПВ от многократных включений (Блок\_повт\_ЧАПВ)  & \centering SGF4 & \centering Не предусмотрено \\ Предусмотрено & \centering -- & \centering \arraybackslash -- \\
\hline
\centering Тип пуска ЧАПВ (Тип\_пуска)  & \centering SGF3 & \centering Внутренний \\ Внешний & \centering -- & \centering \arraybackslash -- \\
\hline
\centering Уровень частоты срабатывания (Fср)  & \centering F> & \centering 49,0...55,0 & \centering Гц & \centering \arraybackslash 0,1 \\
\hline
\centering Разность частот срабатывания и возврата (dFвоз)  & \centering -- & \centering 0,04...1,00 & \centering Гц & \centering \arraybackslash 0,01 \\
\hline
\centering Напряжение срабатывания (Uср)  & \centering U< & \centering 2,0...100,0 & \centering В & \centering \arraybackslash 0,1 \\
\hline
\centering Режим работы функции (Реж\_работы)  & \centering SGF2 & \centering По входу \\ Измеренный & \centering -- & \centering \arraybackslash -- \\
\hline
\centering Время готовности ЧАПВ (Tгот)  & \centering T1 & \centering 5000...240000 & \centering мс & \centering \arraybackslash 10 \\
\hline
\centering Выдержка времени срабатывания (Tср)  & \centering T2 & \centering 0...100000 & \centering мс & \centering \arraybackslash 10 \\
\hline
%===t1
\end{longtable}
\normalsize


%%%%%%%%%%%%%%%%%%%%%%%%%%%%%%%%%%%%%%%%%%%%%%%%%%%%%%%%%% АПВ %%%%%%%%%%%%%%%%%%%%%%%%%%%%%%%%%%%%%%%%%%%%%%%%%%%%%%%%%%%%%%%%%%
\par
В таблице \ref{apv:tbl1} приведены параметры, необходимые для настройки АПВ.

\footnotesize
\begin{longtable}{|>{\centering\arraybackslash}m{5.3cm}|>{\centering\arraybackslash}m{3.3cm}|>{\centering\arraybackslash}m{4.2cm}|>{\centering\arraybackslash}m{1.8cm}|>{\centering\arraybackslash}m{1cm}|}
\caption{Параметры для настройки АПВ\hfill\vspace{-0.5\baselineskip}}\label{apv:tbl1}\\ 
\hline
\rowcolor{gray!20}
Параметр (Параметр на ИЧМ) & Условное обозначение на схеме & Значение/ Диапазон & Единица измерения & Шаг \\ 
\hline
\endfirsthead
\caption*{\hspace{3pt}\emph{Продолжение таблицы \ref{apv:tbl1}\hfill\vspace{-0.5\baselineskip}}} \\ % сделано по ГОСТ 2.105 п.6.8.7
\hline
\rowcolor{gray!20}
Параметр (Параметр на ИЧМ) & Условное обозначение на схеме & Значение/ Диапазон & Единица измерения & Шаг \\ 
%\hline
\endhead
%\hline
\endfoot
%\hline
\endlastfoot
%==+t1*RREC1|LVLINREC
\centering Ввод функции в работу (Ввод\_функции)  & \centering SGF1 & \centering Не предусмотрено \\ Предусмотрено & \centering -- & \centering \arraybackslash -- \\
\hline
\centering Количество циклов АПВ (Число\_циклов)  & \centering SGF4 & \centering 1 \\ 2 & \centering -- & \centering \arraybackslash -- \\
\hline
\centering Режим блокировки второго цикла при ОЗЗ (Блок\_2ц\_ОЗЗ)  & \centering SGF5 & \centering Не предусмотрено \\ Предусмотрено & \centering -- & \centering \arraybackslash -- \\
\hline
\centering Режим контроля напряжения (Реж\_конт\_U)  & \centering SGF2 & \centering Не предусмотрено \\ Оперативный \\ КННш+КОНп \\ КОНш+КННп \\ КННш+КОНп-КОНш+КННп & \centering -- & \centering \arraybackslash -- \\
\hline
\centering Контроль синхронизма и напряжений (Контр\_синхр)  & \centering SGF3 & \centering Не предусмотрено \\ Предусмотрено & \centering -- & \centering \arraybackslash -- \\
\hline
\centering Время готовности для однократного АПВ (Тгот\_1ц)  & \centering T1 & \centering 1000...100000 & \centering мс & \centering \arraybackslash 10 \\
\hline
\centering Время готовности для двукратного АПВ (Тгот\_2ц)  & \centering T6 & \centering 50000...150000 & \centering мс & \centering \arraybackslash 10 \\
\hline
\centering Выдержка времени первого цикла АПВ (Тср)  & \centering T2 & \centering 100...180000 & \centering мс & \centering \arraybackslash 10 \\
\hline
\centering Выдержка времени второго цикла АПВ (Тср\_2ц)  & \centering T5 & \centering 3000...300000 & \centering мс & \centering \arraybackslash 10 \\
\hline
\centering Длительность задержки включения (Тзад\_вкл)  & \centering T3 & \centering 1000...3600000 & \centering мс & \centering \arraybackslash 1000 \\
\hline
%===t1
\end{longtable}
\normalsize



%%%%%%%%%%%%%%%%%%%%%%%%%%%%%%%%%%%%%%%%%%%%%%%%%%%%%%%%%% КННш %%%%%%%%%%%%%%%%%%%%%%%%%%%%%%%%%%%%%%%%%%%%%%%%%%%%%%%%%%%%%%%%%%
\par
В таблице \ref{knnsh:tbl1} приведены параметры, необходимые для настройки КННш.

\footnotesize
\begin{longtable}{|>{\centering\arraybackslash}m{5.3cm}|>{\centering\arraybackslash}m{3.3cm}|>{\centering\arraybackslash}m{4.2cm}|>{\centering\arraybackslash}m{1.8cm}|>{\centering\arraybackslash}m{1cm}|}
\caption{Параметры для настройки КННш\hfill\vspace{-0.5\baselineskip}}\label{knnsh:tbl1}\\ 
\hline
\rowcolor{gray!20}
Параметр (Параметр на ИЧМ) & Условное обозначение на схеме & Значение/ Диапазон & Единица измерения & Шаг \\ 
\hline
\endfirsthead
\caption*{\hspace{3pt}\emph{Продолжение таблицы \ref{knnsh:tbl1}\hfill\vspace{-0.5\baselineskip}}} \\ % сделано по ГОСТ 2.105 п.6.8.7
\hline
\rowcolor{gray!20}
Параметр (Параметр на ИЧМ) & Условное обозначение на схеме & Значение/ Диапазон & Единица измерения & Шаг \\ 
%\hline
\endhead
%\hline
\endfoot
%\hline
\endlastfoot
%==+t1*PTOV1|VOLPRES1
\centering Ввод функции в работу (Ввод\_функции)  & \centering SGF1 & \centering Не предусмотрено \\ Предусмотрено & \centering -- & \centering \arraybackslash -- \\
\hline
\centering Напряжение срабатывания (Uср)  & \centering U> & \centering 50,0...150,0 & \centering В & \centering \arraybackslash 0,1 \\
\hline
\centering Напряжение срабатывания ОП (U2ср)  & \centering U2> & \centering 2,0...50,0 & \centering В & \centering \arraybackslash 0,1 \\
\hline
\centering Режим пуска (Реж\_пуска)  & \centering SGF2 & \centering Измеренный \\ Внешний & \centering -- & \centering \arraybackslash -- \\
\hline
%===t1
\end{longtable}
\normalsize


%%%%%%%%%%%%%%%%%%%%%%%%%%%%%%%%%%%%%%%%%%%%%%%%%%%%%%%%%% КННп %%%%%%%%%%%%%%%%%%%%%%%%%%%%%%%%%%%%%%%%%%%%%%%%%%%%%%%%%%%%%%%%%%
\par
В таблице \ref{knnp:tbl1} приведены параметры, необходимые для настройки КННп.

\footnotesize
\begin{longtable}{|>{\centering\arraybackslash}m{5.3cm}|>{\centering\arraybackslash}m{3.3cm}|>{\centering\arraybackslash}m{4.2cm}|>{\centering\arraybackslash}m{1.8cm}|>{\centering\arraybackslash}m{1cm}|}
\caption{Параметры для настройки КННп\hfill\vspace{-0.5\baselineskip}}\label{knnp:tbl1}\\ 
\hline
\rowcolor{gray!20}
Параметр (Параметр на ИЧМ) & Условное обозначение на схеме & Значение/ Диапазон & Единица измерения & Шаг \\ 
\hline
\endfirsthead
\caption*{\hspace{3pt}\emph{Продолжение таблицы \ref{knnp:tbl1}\hfill\vspace{-0.5\baselineskip}}} \\ % сделано по ГОСТ 2.105 п.6.8.7
\hline
\rowcolor{gray!20}
Параметр (Параметр на ИЧМ) & Условное обозначение на схеме & Значение/ Диапазон & Единица измерения & Шаг \\ 
%\hline
\endhead
%\hline
\endfoot
%\hline
\endlastfoot
%==+t1*PTOV1|VOLPRES2
\centering Ввод функции в работу (Ввод\_функции)  & \centering SGF1 & \centering Не предусмотрено \\ Предусмотрено & \centering -- & \centering \arraybackslash -- \\
\hline
\centering Напряжение срабатывания (Uср)  & \centering U> & \centering 50,0...150,0 & \centering В & \centering \arraybackslash 0,1 \\
\hline
\centering Напряжение срабатывания ОП (U2ср)  & \centering U2> & \centering 2,0...50,0 & \centering В & \centering \arraybackslash 0,1 \\
\hline
\centering Контроль ТН (Контроль\_ТН)  & \centering SGF3 & \centering 3 фазы \\ 2 фазы & \centering -- & \centering \arraybackslash -- \\
\hline
\centering Режим пуска (Реж\_пуска)  & \centering SGF2 & \centering Измеренный \\ Внешний & \centering -- & \centering \arraybackslash -- \\
\hline
%===t1
\end{longtable}
\normalsize


%%%%%%%%%%%%%%%%%%%%%%%%%%%%%%%%%%%%%%%%%%%%%%%%%%%%%%%%%% КОНш %%%%%%%%%%%%%%%%%%%%%%%%%%%%%%%%%%%%%%%%%%%%%%%%%%%%%%%%%%%%%%%%%%
\par
В таблице \ref{konsh:tbl1} приведены параметры, необходимые для настройки КОНш.

\footnotesize
\begin{longtable}{|>{\centering\arraybackslash}m{5.3cm}|>{\centering\arraybackslash}m{3.3cm}|>{\centering\arraybackslash}m{4.2cm}|>{\centering\arraybackslash}m{1.8cm}|>{\centering\arraybackslash}m{1cm}|}
\caption{Параметры для настройки КОНш\hfill\vspace{-0.5\baselineskip}}\label{konsh:tbl1}\\ 
\hline
\rowcolor{gray!20}
Параметр (Параметр на ИЧМ) & Условное обозначение на схеме & Значение/ Диапазон & Единица измерения & Шаг \\ 
\hline
\endfirsthead
\caption*{\hspace{3pt}\emph{Продолжение таблицы \ref{konsh:tbl1}\hfill\vspace{-0.5\baselineskip}}} \\ % сделано по ГОСТ 2.105 п.6.8.7
\hline
\rowcolor{gray!20}
Параметр (Параметр на ИЧМ) & Условное обозначение на схеме & Значение/ Диапазон & Единица измерения & Шаг \\ 
%\hline
\endhead
%\hline
\endfoot
%\hline
\endlastfoot
%==+t1*PTUV1|VOLABS1
\centering Ввод функции в работу (Ввод\_функции)  & \centering SGF1 & \centering Не предусмотрено \\ Предусмотрено & \centering -- & \centering \arraybackslash -- \\
\hline
\centering Напряжение срабатывания (Uср)  & \centering U< & \centering 2,0...100,0 & \centering В & \centering \arraybackslash 0,1 \\
\hline
\centering Напряжение срабатывания ОП (U2ср)  & \centering U2> & \centering 2,0...50,0 & \centering В & \centering \arraybackslash 0,1 \\
\hline
\centering Режим пуска (Реж\_пуска)  & \centering SGF2 & \centering Измеренный \\ Внешний & \centering -- & \centering \arraybackslash -- \\
\hline
%===t1
\end{longtable}
\normalsize


%%%%%%%%%%%%%%%%%%%%%%%%%%%%%%%%%%%%%%%%%%%%%%%%%%%%%%%%%% КОНп %%%%%%%%%%%%%%%%%%%%%%%%%%%%%%%%%%%%%%%%%%%%%%%%%%%%%%%%%%%%%%%%%%
\par
В таблице \ref{konp:tbl1} приведены параметры, необходимые для настройки КОНп.

\footnotesize
\begin{longtable}{|>{\centering\arraybackslash}m{5.3cm}|>{\centering\arraybackslash}m{3.3cm}|>{\centering\arraybackslash}m{4.2cm}|>{\centering\arraybackslash}m{1.8cm}|>{\centering\arraybackslash}m{1cm}|}
\caption{Параметры для настройки КОНп\hfill\vspace{-0.5\baselineskip}}\label{konp:tbl1}\\ 
\hline
\rowcolor{gray!20}
Параметр (Параметр на ИЧМ) & Условное обозначение на схеме & Значение/ Диапазон & Единица измерения & Шаг \\ 
\hline
\endfirsthead
\caption*{\hspace{3pt}\emph{Продолжение таблицы \ref{konp:tbl1}\hfill\vspace{-0.5\baselineskip}}} \\ % сделано по ГОСТ 2.105 п.6.8.7
\hline
\rowcolor{gray!20}
Параметр (Параметр на ИЧМ) & Условное обозначение на схеме & Значение/ Диапазон & Единица измерения & Шаг \\ 
%\hline
\endhead
%\hline
\endfoot
%\hline
\endlastfoot
%==+t1*PTUV1|VOLABS2
\centering Ввод функции в работу (Ввод\_функции)  & \centering SGF1 & \centering Не предусмотрено \\ Предусмотрено & \centering -- & \centering \arraybackslash -- \\
\hline
\centering Напряжение срабатывания (Uср)  & \centering U< & \centering 2,0...100,0 & \centering В & \centering \arraybackslash 0,1 \\
\hline
\centering Напряжение срабатывания ОП (U2ср)  & \centering U2> & \centering 2,0...50,0 & \centering В & \centering \arraybackslash 0,1 \\
\hline
\centering Контроль ТН (Контроль\_ТН)  & \centering SGF3 & \centering 3 фазы \\ 2 фазы & \centering -- & \centering \arraybackslash -- \\
\hline
\centering Режим пуска (Реж\_пуска)  & \centering SGF2 & \centering Измеренный \\ Внешний & \centering -- & \centering \arraybackslash -- \\
\hline
%===t1
\end{longtable}
\normalsize


%%%%%%%%%%%%%%%%%%%%%%%%%%%%%%%%%%%%%%%%%%%%%%%%%%%%%%%%%% КС %%%%%%%%%%%%%%%%%%%%%%%%%%%%%%%%%%%%%%%%%%%%%%%%%%%%%%%%%%%%%%%%%%
\par
В таблице \ref{ks:tbl1} приведены параметры, необходимые для настройки КС.

\footnotesize
\begin{longtable}{|>{\centering\arraybackslash}m{5.3cm}|>{\centering\arraybackslash}m{3.3cm}|>{\centering\arraybackslash}m{4.2cm}|>{\centering\arraybackslash}m{1.8cm}|>{\centering\arraybackslash}m{1cm}|}
\caption{Параметры для настройки КС\hfill\vspace{-0.5\baselineskip}}\label{ks:tbl1}\\ 
\hline
\rowcolor{gray!20}
Параметр (Параметр на ИЧМ) & Условное обозначение на схеме & Значение/ Диапазон & Единица измерения & Шаг \\ 
\hline
\endfirsthead
\caption*{\hspace{3pt}\emph{Продолжение таблицы \ref{ks:tbl1}\hfill\vspace{-0.5\baselineskip}}} \\ % сделано по ГОСТ 2.105 п.6.8.7
\hline
\rowcolor{gray!20}
Параметр (Параметр на ИЧМ) & Условное обозначение на схеме & Значение/ Диапазон & Единица измерения & Шаг \\ 
%\hline
\endhead
%\hline
\endfoot
%\hline
\endlastfoot
%==+t1*LVRSYN1|LVSYN
\centering Ввод функции в работу (Ввод\_функции)  & \centering SGF1 & \centering Не предусмотрено \\ Предусмотрено & \centering -- & \centering \arraybackslash -- \\
\hline
\centering Синхронизируемое напряжение (Синхрон\_напряж)  & \centering SGF2 & \centering UAB \\ UBC \\ UCA & \centering -- & \centering \arraybackslash -- \\
\hline
\centering Допустимая разность модуля векторов напряжений (dUср)  & \centering dU< & \centering 5,0...50,0 & \centering В & \centering \arraybackslash 0,1 \\
\hline
\centering Допустимая разность частот для ожидания синхронизма (dFср)  & \centering dFос< & \centering 0,05...0,50 & \centering Гц & \centering \arraybackslash 0,01 \\
\hline
\centering Разность частот срабатывания и возврата ИО частоты ожидания синхронизма (dFвоз)  & \centering -- & \centering 0,04...0,10 & \centering Гц & \centering \arraybackslash 0,01 \\
\hline
\centering Максимальная разность частот для улавливания синхронизма (dFмакс\_ср)  & \centering dFмакс< & \centering 0,05...2,00 & \centering Гц & \centering \arraybackslash 0,01 \\
\hline
\centering Разность частот срабатывания и возврата ИО частоты улавливания синхронизма (dFмакс\_воз)  & \centering -- & \centering 0,04...0,10 & \centering Гц & \centering \arraybackslash 0,01 \\
\hline
\centering Максимальная скорость изменения частоты (dFdtср)  & \centering l dF/dt l> & \centering 0,1...15,0 & \centering Гц/с & \centering \arraybackslash 0,1 \\
\hline
\centering Коэффициент возврата по скорости изменения частоты (Kвоз)  & \centering -- & \centering 0,20...0,99 & \centering -- & \centering \arraybackslash 0,01 \\
\hline
\centering Допустимая разность фаз векторов напряжений (dФср)  & \centering dФос< & \centering 10,0...90,0 & \centering град. & \centering \arraybackslash 0,1 \\
\hline
\centering Угол сдвига (Ф\_сдвига)  & \centering Ф\_сдвига & \centering 0,0...359,9 & \centering град. & \centering \arraybackslash 0,1 \\
\hline
\centering Улавливание синхронизма (Улавл\_синхр)  & \centering SGF3 & \centering Не предусмотрено \\ Предусмотрено & \centering -- & \centering \arraybackslash -- \\
\hline
\centering Максимально допустимый угол включения выключателя (dФус)  & \centering dФус< & \centering 10,0...90,0 & \centering град. & \centering \arraybackslash 0,1 \\
\hline
\centering Время включения выключателя (Tопер)  & \centering Tопер & \centering 10...2000 & \centering мс & \centering \arraybackslash 10 \\
\hline
%===t1
\end{longtable}
\normalsize


%%%%%%%%%%%%%%%%%%%%%%%%%%%%%%%%%%%%%%%%%%%%%%%%%%%%%%%%%% ЛОВ %%%%%%%%%%%%%%%%%%%%%%%%%%%%%%%%%%%%%%%%%%%%%%%%%%%%%%%%%%%%%%%%%%
\par
В таблице \ref{lorz:tbl1} приведены параметры, необходимые для настройки ЛОВ.

\footnotesize
\begin{longtable}{|>{\centering\arraybackslash}m{5.3cm}|>{\centering\arraybackslash}m{3.3cm}|>{\centering\arraybackslash}m{4.2cm}|>{\centering\arraybackslash}m{1.8cm}|>{\centering\arraybackslash}m{1cm}|}
\caption{Параметры для настройки ЛОВ\hfill\vspace{-0.5\baselineskip}}\label{lorz:tbl1}\\ 
\hline
\rowcolor{gray!20}
Параметр (Параметр на ИЧМ) & Условное обозначение на схеме & Значение/ Диапазон & Единица измерения & Шаг \\ 
\hline
\endfirsthead
\caption*{\hspace{3pt}\emph{Продолжение таблицы \ref{lorz:tbl1}\hfill\vspace{-0.5\baselineskip}}} \\ % сделано по ГОСТ 2.105 п.6.8.7
\hline
\rowcolor{gray!20}
Параметр (Параметр на ИЧМ) & Условное обозначение на схеме & Значение/ Диапазон & Единица измерения & Шаг \\ 
%\hline
\endhead
%\hline
\endfoot
%\hline
\endlastfoot
%==+t1*OFFPTRC1|LVLINONOFFLGC
\multicolumn{5}{|c|}{ ЛО РЗ } \\ \hline 
\centering Ввод функции в работу (Ввод\_функции)  & \centering SGF1 & \centering Не предусмотрено \\ Предусмотрено & \centering -- & \centering \arraybackslash -- \\
\hline
\centering Длительность импульса (Tимп)  & \centering Т1 & \centering 50...60000 & \centering мс & \centering \arraybackslash 10 \\
\hline
\multicolumn{5}{|c|}{ ЛО ПА } \\ \hline 
\centering Ввод функции в работу (Ввод\_функции)  & \centering SGF1 & \centering Не предусмотрено \\ Предусмотрено & \centering -- & \centering \arraybackslash -- \\
\hline
\centering Длительность импульса (Tимп)  & \centering Т1 & \centering 50...60000 & \centering мс & \centering \arraybackslash 10 \\
\hline
\multicolumn{5}{|c|}{ ЛВ ПА } \\ \hline 
\centering Ввод функции в работу (Ввод\_функции)  & \centering SGF1 & \centering Не предусмотрено \\ Предусмотрено & \centering -- & \centering \arraybackslash -- \\
\hline
\centering Тип пуска (Тип\_пуска)  & \centering SGF2 & \centering Внешний \\ Внутренний & \centering -- & \centering \arraybackslash -- \\
\hline
\centering Выдержка времени срабатывания (Тср)  & \centering Т1 & \centering 0...10000 & \centering мс & \centering \arraybackslash 10 \\
\hline
\multicolumn{5}{|c|}{ ЗАПВ } \\ \hline 
\centering Режим запрета АПВ при срабатывании ДЗфф 1 ступени (ЗапАПВ\_ДЗфф\_1ст)  & \centering SGF1 & \centering Не предусмотрено \\ Предусмотрено & \centering -- & \centering \arraybackslash -- \\
\hline
\centering Режим запрета АПВ при срабатывании ДЗфф 2 ступени (ЗапАПВ\_ДЗфф\_2ст)  & \centering SGF2 & \centering Не предусмотрено \\ Предусмотрено & \centering -- & \centering \arraybackslash -- \\
\hline
\centering Режим запрета АПВ при срабатывании ДЗфф 3 ступени (ЗапАПВ\_ДЗфф\_3ст)  & \centering SGF3 & \centering Не предусмотрено \\ Предусмотрено & \centering -- & \centering \arraybackslash -- \\
\hline
\centering Режим запрета АПВ при срабатывании ДЗфз 1 ступени (ЗапАПВ\_ДЗфз\_1ст)  & \centering SGF4 & \centering Не предусмотрено \\ Предусмотрено & \centering -- & \centering \arraybackslash -- \\
\hline
\centering Режим запрета АПВ при срабатывании ДЗфз 2 ступени (ЗапАПВ\_ДЗфз\_2ст)  & \centering SGF5 & \centering Не предусмотрено \\ Предусмотрено & \centering -- & \centering \arraybackslash -- \\
\hline
\centering Режим запрета АПВ при срабатывании НЗОЗЗ (ЗапАПВ\_НЗОЗЗ)  & \centering SGF6 & \centering Не предусмотрено \\ Предусмотрено & \centering -- & \centering \arraybackslash -- \\
\hline
\centering Режим запрета АПВ при срабатывании ТЗНП 1 ступени (ЗапАПВ\_ТЗНП\_1ст)  & \centering SGF7 & \centering Не предусмотрено \\ Предусмотрено & \centering -- & \centering \arraybackslash -- \\
\hline
\centering Режим запрета АПВ при срабатывании ТЗНП 2 ступени (ЗапАПВ\_ТЗНП\_2ст)  & \centering SGF8 & \centering Не предусмотрено \\ Предусмотрено & \centering -- & \centering \arraybackslash -- \\
\hline
\centering Режим запрета АПВ при срабатывании ТО (ЗапАПВ\_ТО)  & \centering SGF9 & \centering Не предусмотрено \\ Предусмотрено & \centering -- & \centering \arraybackslash -- \\
\hline
\centering Режим запрета АПВ при срабатывании МТЗ 1 ступени (ЗапАПВ\_МТЗ\_1ст)  & \centering SGF10 & \centering Не предусмотрено \\ Предусмотрено & \centering -- & \centering \arraybackslash -- \\
\hline
\centering Режим запрета АПВ при срабатывании МТЗ 2 ступени (ЗапАПВ\_МТЗ\_2ст)  & \centering SGF11 & \centering Не предусмотрено \\ Предусмотрено & \centering -- & \centering \arraybackslash -- \\
\hline
\centering Режим запрета АПВ при срабатывании МТЗ 3 ступени (ЗапАПВ\_МТЗ\_3ст)  & \centering SGF12 & \centering Не предусмотрено \\ Предусмотрено & \centering -- & \centering \arraybackslash -- \\
\hline
\centering Режим запрета АПВ при срабатывании ЗОП (ЗапАПВ\_ЗОП)  & \centering SGF13 & \centering Не предусмотрено \\ Предусмотрено & \centering -- & \centering \arraybackslash -- \\
\hline
\centering Режим запрета АПВ при срабатывании АУ (ЗапАПВ\_АУ)  & \centering SGF14 & \centering Не предусмотрено \\ Предусмотрено & \centering -- & \centering \arraybackslash -- \\
\hline
\centering Режим запрета АПВ при срабатывании ЗП (ЗапАПВ\_ЗП)  & \centering SGF15 & \centering Не предусмотрено \\ Предусмотрено & \centering -- & \centering \arraybackslash -- \\
\hline
\centering Режим запрета АПВ при срабатывании ГСОЗЗ (ЗапАПВ\_ГСОЗЗ)  & \centering SGF16 & \centering Не предусмотрено \\ Предусмотрено & \centering -- & \centering \arraybackslash -- \\
\hline
\centering Режим запрета АПВ при срабатывании ГЗоткл (ЗапАПВ\_Сраб\_ГЗоткл)  & \centering SGF17 & \centering Не предусмотрено \\ Предусмотрено & \centering -- & \centering \arraybackslash -- \\
\hline
\centering Режим запрета АПВ при срабатывании ГЗсигн (ЗапАПВ\_Сраб\_ГЗсигн)  & \centering SGF18 & \centering Не предусмотрено \\ Предусмотрено & \centering -- & \centering \arraybackslash -- \\
\hline
\centering Режим запрета АПВ при срабатывании ДТм (ЗапАПВ\_Сраб\_ДТм)  & \centering SGF19 & \centering Не предусмотрено \\ Предусмотрено & \centering -- & \centering \arraybackslash -- \\
\hline
\centering Режим запрета АПВ при срабатывании ДТо (ЗапАПВ\_Сраб\_ДТо)  & \centering SGF20 & \centering Не предусмотрено \\ Предусмотрено & \centering -- & \centering \arraybackslash -- \\
\hline
\centering Режим запрета АПВ при срабатывании ЛО РД (ЗапАПВ\_Сраб\_ЛО\_РД)  & \centering SGF21 & \centering Не предусмотрено \\ Предусмотрено & \centering -- & \centering \arraybackslash -- \\
\hline
\centering Режим запрета АПВ при срабатывании ОУ ДЗ (ЗапАПВ\_ОУ\_ДЗ)  & \centering SGF22 & \centering Не предусмотрено \\ Предусмотрено & \centering -- & \centering \arraybackslash -- \\
\hline
\centering Режим запрета АПВ при срабатывании ОУ ТЗНП (ЗапАПВ\_ОУ\_ТЗНП)  & \centering SGF23 & \centering Не предусмотрено \\ Предусмотрено & \centering -- & \centering \arraybackslash -- \\
\hline
\centering Режим запрета АПВ при срабатывании ОУ МТЗ (ЗапАПВ\_ОУ\_МТЗ)  & \centering SGF24 & \centering Не предусмотрено \\ Предусмотрено & \centering -- & \centering \arraybackslash -- \\
\hline
\centering Режим запрета АПВ при срабатывании ЗПЧ 1 ступени (ЗапАПВ\_ЗПЧ\_1ст)  & \centering SGF25 & \centering Не предусмотрено \\ Предусмотрено & \centering -- & \centering \arraybackslash -- \\
\hline
\centering Режим запрета АПВ при срабатывании ЗПЧ 2 ступени (ЗапАПВ\_ЗПЧ\_2ст)  & \centering SGF26 & \centering Не предусмотрено \\ Предусмотрено & \centering -- & \centering \arraybackslash -- \\
\hline
\centering Режим запрета АПВ при срабатывании ЗИЧ 1 ступени (ЗапАПВ\_ЗИЧ\_1ст)  & \centering SGF27 & \centering Не предусмотрено \\ Предусмотрено & \centering -- & \centering \arraybackslash -- \\
\hline
\centering Режим запрета АПВ при срабатывании ЗИЧ 2 ступени (ЗапАПВ\_ЗИЧ\_2ст)  & \centering SGF28 & \centering Не предусмотрено \\ Предусмотрено & \centering -- & \centering \arraybackslash -- \\
\hline
\centering Режим запрета АПВ при срабатывании ЗСЧ 1 ступени (ЗапАПВ\_ЗСЧ\_1ст)  & \centering SGF29 & \centering Не предусмотрено \\ Предусмотрено & \centering -- & \centering \arraybackslash -- \\
\hline
\centering Режим запрета АПВ при срабатывании ЗСЧ 2 ступени (ЗапАПВ\_ЗСЧ\_2ст)  & \centering SGF30 & \centering Не предусмотрено \\ Предусмотрено & \centering -- & \centering \arraybackslash -- \\
\hline
\centering Режим запрета АПВ при наличии отключения от АЧР (ЗапАПВ\_АЧР)  & \centering SGF31 & \centering Не предусмотрено \\ Предусмотрено & \centering -- & \centering \arraybackslash -- \\
\hline
\centering Режим запрета АПВ при срабатывании ЗДЗ (ЗапАПВ\_ЗДЗ)  & \centering SGF32 & \centering Не предусмотрено \\ Предусмотрено & \centering -- & \centering \arraybackslash -- \\
\hline
\centering Режим запрета АПВ при срабатывании ЗМН 1 ступени (ЗапАПВ\_ЗМН\_1ст)  & \centering SGF33 & \centering Не предусмотрено \\ Предусмотрено & \centering -- & \centering \arraybackslash -- \\
\hline
\centering Режим запрета АПВ при срабатывании ЗМН 2 ступени (ЗапАПВ\_ЗМН\_2ст)  & \centering SGF34 & \centering Не предусмотрено \\ Предусмотрено & \centering -- & \centering \arraybackslash -- \\
\hline
\centering Режим запрета АПВ при срабатывании ЗПН 1 ступени (ЗапАПВ\_ЗПН\_1ст)  & \centering SGF35 & \centering Не предусмотрено \\ Предусмотрено & \centering -- & \centering \arraybackslash -- \\
\hline
\centering Режим запрета АПВ при срабатывании ЗПН 2 ступени (ЗапАПВ\_ЗПН\_2ст)  & \centering SGF36 & \centering Не предусмотрено \\ Предусмотрено & \centering -- & \centering \arraybackslash -- \\
\hline
\centering Режим запрета АПВ при срабатывании УРОВ «на себя» (ЗапАПВ\_УРОВ\_на\_себя)  & \centering SGF37 & \centering Не предусмотрено \\ Предусмотрено & \centering -- & \centering \arraybackslash -- \\
\hline
\centering Режим запрета АПВ при наличии отключения от ЛО ПА (ЗапАПВ\_Откл\_ЛО\_ПА)  & \centering SGF38 & \centering Не предусмотрено \\ Предусмотрено & \centering -- & \centering \arraybackslash -- \\
\hline
\centering Режим запрета АПВ при наличии отключения внешнего 1 (ЗапАПВ\_Откл\_внеш\_1)  & \centering SGF39 & \centering Не предусмотрено \\ Предусмотрено & \centering -- & \centering \arraybackslash -- \\
\hline
\centering Режим запрета АПВ при наличии отключения внешнего 2 (ЗапАПВ\_Откл\_внеш\_2)  & \centering SGF40 & \centering Не предусмотрено \\ Предусмотрено & \centering -- & \centering \arraybackslash -- \\
\hline
\centering Режим запрета АПВ при наличии отключения внешнего 3 (ЗапАПВ\_Откл\_внеш\_3)  & \centering SGF41 & \centering Не предусмотрено \\ Предусмотрено & \centering -- & \centering \arraybackslash -- \\
\hline
\centering Режим запрета АПВ при наличии отключения внешнего 4 (ЗапАПВ\_Откл\_внеш\_4)  & \centering SGF42 & \centering Не предусмотрено \\ Предусмотрено & \centering -- & \centering \arraybackslash -- \\
\hline
%===t1
\end{longtable}
\normalsize


%%%%%%%%%%%%%%%%%%%%%%%%%%%%%%%%%%%%%%%%%%%%%%%%%%%%%%%%%% КП %%%%%%%%%%%%%%%%%%%%%%%%%%%%%%%%%%%%%%%%%%%%%%%%%%%%%%%%%%%%%%%%%%
\par
В таблице \ref{kp:tbl_gen} приведены параметры, необходимые для настройки КП. 

\footnotesize
\begin{longtable}{|>{\centering\arraybackslash}m{5.3cm}|>{\centering\arraybackslash}m{3.3cm}|>{\centering\arraybackslash}m{4.2cm}|>{\centering\arraybackslash}m{1.8cm}|>{\centering\arraybackslash}m{1cm}|}
\caption{Параметры для настройки КП\hfill\vspace{-0.5\baselineskip}}\label{kp:tbl_gen}\\ 
\hline
\rowcolor{gray!20}
Параметр (Параметр на ИЧМ) & Условное обозначение на схеме & Значение/ Диапазон & Единица измерения & Шаг \\ 
\hline
\endfirsthead
\caption*{\hspace{3pt}\emph{Продолжение таблицы \ref{kp:tbl_gen}\hfill\vspace{-0.5\baselineskip}}} \\ % сделано по ГОСТ 2.105 п.6.8.7
\hline
\rowcolor{gray!20}
Параметр (Параметр на ИЧМ) & Условное обозначение на схеме & Значение/ Диапазон & Единица измерения & Шаг \\ 
%\hline
\endhead
%\hline
\endfoot
%\hline
\endlastfoot
%==+t1*LLN0|SWCTRL
\multicolumn{5}{|c|}{ УВ } \\ \hline 
\centering Ввод функции в работу (Ввод\_функции)  & \centering SGF1 & \centering Не предусмотрено \\ Предусмотрено & \centering -- & \centering \arraybackslash -- \\
\hline
\centering Включение с контролем синхронизма (Контроль\_синхр)  & \centering SGF2 & \centering Не предусмотрено \\ Предусмотрено & \centering -- & \centering \arraybackslash -- \\
\hline
\centering Блокировка включения от аварийного отключения В (Бл\_вкл\_от\_авар\_откл)  & \centering SGF3 & \centering Не предусмотрено \\ Предусмотрено & \centering -- & \centering \arraybackslash -- \\
\hline
\centering Допустимое время переключения выключателя (Tперекл)  & \centering T1 & \centering 10...10000 & \centering мс & \centering \arraybackslash 10 \\
\hline
\centering Выдержка времени подавления выдачи положения "Не определено" (Tблк)  & \centering T2 & \centering 10...10000 & \centering мс & \centering \arraybackslash 10 \\
\hline
\centering Длительность задержки включения с контролем синхронизма (Тзад\_вкл\_с\_КС)  & \centering T3 & \centering 50...60000 & \centering мс & \centering \arraybackslash 10 \\
\hline
\multicolumn{5}{|c|}{ УВЭ } \\ \hline 
\centering Ввод функции в работу (Ввод\_функции)  & \centering SGF1 & \centering Не предусмотрено \\ Предусмотрено & \centering -- & \centering \arraybackslash -- \\
\hline
\centering Допустимое время переключения (Tперекл)  & \centering T1 & \centering 50...100000 & \centering мс & \centering \arraybackslash 10 \\
\hline
\centering Выдержка времени подавления выдачи положения "Не определено" (Тблк)  & \centering T2 & \centering 50...100000 & \centering мс & \centering \arraybackslash 10 \\
\hline
\multicolumn{5}{|c|}{ УЗН } \\ \hline 
\centering Ввод функции в работу (Ввод\_функции)  & \centering SGF1 & \centering Не предусмотрено \\ Предусмотрено & \centering -- & \centering \arraybackslash -- \\
\hline
\centering Допустимое время переключения (Tперекл)  & \centering T1 & \centering 50...100000 & \centering мс & \centering \arraybackslash 10 \\
\hline
\centering Выдержка времени подавления выдачи положения "Не определено" (Тблк)  & \centering T2 & \centering 50...100000 & \centering мс & \centering \arraybackslash 10 \\
\hline
\multicolumn{5}{|c|}{ ОБР В } \\ \hline 
\centering Ввод ОБР В в работу (ОБР\_В)  & \centering SGF1 & \centering Не предусмотрено \\ Предусмотрено & \centering -- & \centering \arraybackslash -- \\
\hline
\centering Контроль ЗН (Контр\_ЗН)  & \centering SGF2 & \centering Не предусмотрено \\ Предусмотрено & \centering -- & \centering \arraybackslash -- \\
\hline
\centering Контроль КА1 (Контр\_KA1)  & \centering SGF3 & \centering Не предусмотрено \\ Предусмотрено & \centering -- & \centering \arraybackslash -- \\
\hline
\centering Контроль КА2 (Контр\_KA2)  & \centering SGF4 & \centering Не предусмотрено \\ Предусмотрено & \centering -- & \centering \arraybackslash -- \\
\hline
\multicolumn{5}{|c|}{ ОБР ВЭ } \\ \hline 
\centering Контроль КА1 (Контр\_КА1)  & \centering SGF1 & \centering Не предусмотрено \\ Предусмотрено & \centering -- & \centering \arraybackslash -- \\
\hline
\centering Контроль КА2 (Контр\_КА2)  & \centering SGF2 & \centering Не предусмотрено \\ Предусмотрено & \centering -- & \centering \arraybackslash -- \\
\hline
\multicolumn{5}{|c|}{ ОБР ЗН } \\ \hline 
\centering Контроль KA3 (Контр\_KA3)  & \centering SGF1 & \centering Не предусмотрено \\ Предусмотрено & \centering -- & \centering \arraybackslash -- \\
\hline
\centering Контроль KA4 (Контр\_KA4)  & \centering SGF2 & \centering Не предусмотрено \\ Предусмотрено & \centering -- & \centering \arraybackslash -- \\
\hline
\multicolumn{5}{|c|}{ КП } \\ \hline 
\centering Ввод функции в работу (Ввод\_функции)  & \centering SGF1 & \centering Не предусмотрено \\ Предусмотрено & \centering -- & \centering \arraybackslash -- \\
\hline
\centering Ввод ОБР КА в работу (ОБР\_КА)  & \centering SG1 & \centering Не предусмотрено \\ Предусмотрено & \centering -- & \centering \arraybackslash -- \\
\hline
%===t1
\end{longtable}
\normalsize


%%%%%%%%%%%%%%%%%%%%%%%%%%%%%%%%%%%%%%%%%%%%%%%%%%%%%%%%%% УВ %%%%%%%%%%%%%%%%%%%%%%%%%%%%%%%%%%%%%%%%%%%%%%%%%%%%%%%%%%%%%%%%%%
\par
В таблице \ref{uv:tbl_gen} приведены параметры, необходимые для настройки УВ. 

\footnotesize
\begin{longtable}{|>{\centering\arraybackslash}m{5.3cm}|>{\centering\arraybackslash}m{3.3cm}|>{\centering\arraybackslash}m{4.2cm}|>{\centering\arraybackslash}m{1.8cm}|>{\centering\arraybackslash}m{1cm}|}
\caption{Параметры для настройки УВ\hfill\vspace{-0.5\baselineskip}}\label{uv:tbl_gen}\\ 
\hline
\rowcolor{gray!20}
Параметр (Параметр на ИЧМ) & Условное обозначение на схеме & Значение/ Диапазон & Единица измерения & Шаг \\ 
\hline
\endfirsthead
\caption*{\hspace{3pt}\emph{Продолжение таблицы \ref{uv:tbl_gen}\hfill\vspace{-0.5\baselineskip}}} \\ % сделано по ГОСТ 2.105 п.6.8.7
\hline
\rowcolor{gray!20}
Параметр (Параметр на ИЧМ) & Условное обозначение на схеме & Значение/ Диапазон & Единица измерения & Шаг \\ 
%\hline
\endhead
%\hline
\endfoot
%\hline
\endlastfoot
%==+t1*CBCSWI1|HVBCTRL
\centering Ввод функции в работу (Ввод\_функции)  & \centering SGF1 & \centering Не предусмотрено \\ Предусмотрено & \centering -- & \centering \arraybackslash -- \\
\hline
\centering Включение с контролем синхронизма (Контроль\_синхр)  & \centering SGF2 & \centering Не предусмотрено \\ Предусмотрено & \centering -- & \centering \arraybackslash -- \\
\hline
\centering Длительность задержки включения с контролем синхронизма (Тзад\_вкл\_с\_КС)  & \centering T1 & \centering 50...60000 & \centering мс & \centering \arraybackslash 10 \\
\hline
%===t1
\end{longtable}
\normalsize

%%%%%%%%%%%%%%%%%%%%%%%%%%%%%%%%%%%%%%%%%%%%%%%%%%%%%%%%%% КА %%%%%%%%%%%%%%%%%%%%%%%%%%%%%%%%%%%%%%%%%%%%%%%%%%%%%%%%%%%%%%%%%%
\par
В таблице \ref{ka:tbl_gen} приведены параметры, необходимые для настройки КА. 

\footnotesize
\begin{longtable}{|>{\centering\arraybackslash}m{5.3cm}|>{\centering\arraybackslash}m{3.3cm}|>{\centering\arraybackslash}m{4.2cm}|>{\centering\arraybackslash}m{1.8cm}|>{\centering\arraybackslash}m{1cm}|}
\caption{Параметры для настройки КА\hfill\vspace{-0.5\baselineskip}}\label{ka:tbl_gen}\\ 
\hline
\rowcolor{gray!20}
Параметр (Параметр на ИЧМ) & Условное обозначение на схеме & Значение/ Диапазон & Единица измерения & Шаг \\ 
\hline
\endfirsthead
%\caption{\emph{Продолжение\hfill\vspace{-0.5\baselineskip}}} \\ % Отступ обозначения таблицы от таблицы и выравнивание надписи слева
\caption*{\hspace{3pt}\emph{Продолжение таблицы \ref{ka:tbl_gen}\hfill\vspace{-0.5\baselineskip}}} \\ % сделано по ГОСТ 2.105 п.6.8.7
\hline
\rowcolor{gray!20}
Параметр (Параметр на ИЧМ) & Условное обозначение на схеме & Значение/ Диапазон & Единица измерения & Шаг \\ 
%\hline
\endhead
%\hline
\endfoot
%\hline
\endlastfoot
%==+t1*LLN0|SwitchDevice
\multicolumn{5}{|c|}{ В } \\ \hline 
\centering Ввод функции в работу (Ввод\_функции)  & \centering SGF1 & \centering Не предусмотрено \\ Предусмотрено & \centering -- & \centering \arraybackslash -- \\
\hline
\centering Режим отключения выключателя (Режим\_откл)  & \centering SGF2 & \centering Длительный \\ Импульсный & \centering -- & \centering \arraybackslash -- \\
\hline
\centering Контроль работы ЭМО (Контр\_раб\_ЭМО)  & \centering SGF3 & \centering Не предусмотрено \\ Предусмотрено & \centering -- & \centering \arraybackslash -- \\
\hline
\centering Режим включения выключателя (Режим\_вкл)  & \centering SGF4 & \centering Длительный \\ Импульсный & \centering -- & \centering \arraybackslash -- \\
\hline
\centering Контроль работы ЭМВ (Контр\_раб\_ЭМВ)  & \centering SGF5 & \centering Не предусмотрено \\ Предусмотрено & \centering -- & \centering \arraybackslash -- \\
\hline
\centering Длительность формирования сигнала о неисправном положении В (Tнеиспр\_В)  & \centering T1 & \centering 10...10000 & \centering мс & \centering \arraybackslash 10 \\
\hline
\centering Длительность формирования команды "Отключить В (реле)" (Тимп\_откл)  & \centering T2 & \centering 50...10000 & \centering мс & \centering \arraybackslash 10 \\
\hline
\centering Длительность формирования команды "Включить В (реле)" (Тимп\_вкл)  & \centering T3 & \centering 50...10000 & \centering мс & \centering \arraybackslash 10 \\
\hline
\centering Время для исключения "опрокидывания" В (Тпрод\_вкл)  & \centering T4 & \centering 0...1000 & \centering мс & \centering \arraybackslash 10 \\
\hline
\multicolumn{5}{|c|}{ ЗН } \\ \hline 
\centering Ввод функции в работу (Ввод\_функции)  & \centering SGF1 & \centering Не предусмотрено \\ Предусмотрено & \centering -- & \centering \arraybackslash -- \\
\hline
\centering Режим отключения ЗН (Режим\_откл)  & \centering SGF2 & \centering Длительный \\ Импульсный & \centering -- & \centering \arraybackslash -- \\
\hline
\centering Режим включения ЗН (Режим\_вкл)  & \centering SGF3 & \centering Длительный \\ Импульсный & \centering -- & \centering \arraybackslash -- \\
\hline
\centering Длительность формирования сигнала о неисправном положении ЗН (Tнеиспр\_ЗН)  & \centering T1 & \centering 10...100000 & \centering мс & \centering \arraybackslash 10 \\
\hline
\centering Длительность формирования команды "Отключить ЗН (реле)" (Тимп\_откл)  & \centering T2 & \centering 100...100000 & \centering мс & \centering \arraybackslash 10 \\
\hline
\centering Длительность формирования команды "Включить ЗН (реле)" (Тимп\_вкл)  & \centering T3 & \centering 100...100000 & \centering мс & \centering \arraybackslash 10 \\
\hline
\multicolumn{5}{|c|}{ ВЭ } \\ \hline 
\centering Ввод функции в работу (Ввод\_функции)  & \centering SGF1 & \centering Не предусмотрено \\ Предусмотрено & \centering -- & \centering \arraybackslash -- \\
\hline
\centering Режим работы (выкатить) ВЭ (Режим\_выкат)  & \centering SGF2 & \centering Длительный \\ Импульсный & \centering -- & \centering \arraybackslash -- \\
\hline
\centering Режим работы (вкатить) ВЭ (Режим\_вкат)  & \centering SGF3 & \centering Длительный \\ Импульсный & \centering -- & \centering \arraybackslash -- \\
\hline
\centering Контроль ремонтного положения (Контр\_рем\_полож)  & \centering SGF4 & \centering Не предусмотрено \\ Предусмотрено & \centering -- & \centering \arraybackslash -- \\
\hline
\centering Длительность формирования сигнала о неисправном положении ВЭ (Tнеиспр\_ВЭ)  & \centering T1 & \centering 10...100000 & \centering мс & \centering \arraybackslash 10 \\
\hline
\centering Длительность формирования команды "Выкатить ВЭ (реле)" (Тимп\_выкат)  & \centering T2 & \centering 100...100000 & \centering мс & \centering \arraybackslash 10 \\
\hline
\centering Длительность формирования команды "Вкатить ВЭ (реле)" (Тимп\_вкат)  & \centering T3 & \centering 100...100000 & \centering мс & \centering \arraybackslash 10 \\
\hline
\multicolumn{5}{|c|}{ КА } \\ \hline 
\centering Ввод функции в работу (Ввод\_функции)  & \centering SGF1 & \centering Не предусмотрено \\ Предусмотрено & \centering -- & \centering \arraybackslash -- \\
\hline
%===t1
\end{longtable}
\normalsize


%%%%%%%%%%%%%%%%%%%%%%%%%%%%%%%%%%%%%%%%%%%%%%%%%%%%%%%%%% КСВ %%%%%%%%%%%%%%%%%%%%%%%%%%%%%%%%%%%%%%%%%%%%%%%%%%%%%%%%%%%%%%%%%%
\par
В таблице \ref{ksv:tbl1} приведены параметры, необходимые для настройки КСВ.

\footnotesize
\begin{longtable}{|>{\centering\arraybackslash}m{5.3cm}|>{\centering\arraybackslash}m{3.3cm}|>{\centering\arraybackslash}m{4.2cm}|>{\centering\arraybackslash}m{1.8cm}|>{\centering\arraybackslash}m{1cm}|}
\caption{Параметры для настройки КСВ\hfill\vspace{-0.5\baselineskip}}\label{ksv:tbl1}\\ 
\hline
\rowcolor{gray!20}
Параметр (Параметр на ИЧМ) & Условное обозначение на схеме & Значение/ Диапазон & Единица измерения & Шаг \\ 
\hline
\endfirsthead
\caption*{\hspace{3pt}\emph{Продолжение таблицы \ref{ksv:tbl1}\hfill\vspace{-0.5\baselineskip}}} \\ % сделано по ГОСТ 2.105 п.6.8.7
\hline
\rowcolor{gray!20}
Параметр (Параметр на ИЧМ) & Условное обозначение на схеме & Значение/ Диапазон & Единица измерения & Шаг \\ 
%\hline
\endhead
%\hline
\endfoot
%\hline
\endlastfoot
%==+t1*RCBF1|LVCBSUP
\centering Ввод функции в работу (Ввод\_функции)  & \centering SGF1 & \centering Не предусмотрено \\ Предусмотрено & \centering -- & \centering \arraybackslash -- \\
\hline
\centering Контроль ОТ цепей ЭМВ, ЭМО1 и ЭМО2 (Контроль\_ОТ\_ЭМ)  & \centering SGF5 & \centering Не предусмотрено \\ ЭМВ и ЭМО1 \\ ЭМВ, ЭМО1 и ЭМО2 & \centering -- & \centering \arraybackslash -- \\
\hline
\centering Блокировка управления при снижении уровня изоляции В (Блок\_упр\_КИ\_В)  & \centering SGF8 & \centering Не предусмотрено \\ От аварийного \\ От аварийного и низкого & \centering -- & \centering \arraybackslash -- \\
\hline
\centering Блокировка включения при низком уровне изоляции В (Блок\_вкл\_низ\_изол\_В)  & \centering SGF2 & \centering Не предусмотрено \\ Предусмотрено & \centering -- & \centering \arraybackslash -- \\
\hline
\centering Блокировка включения при неисправности положения В (Блок\_вкл\_полож\_В)  & \centering SGF3 & \centering Не предусмотрено \\ Предусмотрено & \centering -- & \centering \arraybackslash -- \\
\hline
\centering Блокировка включения при превышении ресурса В (Блок\_вкл\_ресурса\_В)  & \centering SGF4 & \centering Не предусмотрено \\ Предусмотрено & \centering -- & \centering \arraybackslash -- \\
\hline
\centering Контроль ЭМВ, ЭМО1 и ЭМО2 при формировании неисправности цепей ЭМУ (Контроль\_ЭМ)  & \centering SGF6 & \centering Не предусмотрено \\ ЭМВ и ЭМО1 \\ ЭМВ, ЭМО1 и ЭМО2 & \centering -- & \centering \arraybackslash -- \\
\hline
\centering Разрешение сброса "РФК" от кнопки (Контроль\_кнопки)  & \centering SGF7 & \centering Не предусмотрено \\ Предусмотрено & \centering -- & \centering \arraybackslash -- \\
\hline
\centering Выдержка времени ожидания готовности привода выключателя (Tгот\_в)  & \centering T1 & \centering 0...100000 & \centering мс & \centering \arraybackslash 10 \\
\hline
\centering Выдержка времени срабатывания неисправности цепей ЭМУ (Tдлит\_неисп\_ЭМУ)  & \centering T2 & \centering 0...10000 & \centering мс & \centering \arraybackslash 10 \\
\hline
\centering Выдержка времени срабатывания защит электромагнитов (Tдлит\_ЭМ)  & \centering T3 & \centering 0...10000 & \centering мс & \centering \arraybackslash 10 \\
\hline
%===t1
\end{longtable}
\normalsize



%%%%%%%%%%%%%%%%%%%%%%%%%%%%%%%%%%%%%%%%%%%%%%%%%%%%%%%%%% КРВ %%%%%%%%%%%%%%%%%%%%%%%%%%%%%%%%%%%%%%%%%%%%%%%%%%%%%%%%%%%%%%%%%%
\par
В таблице \ref{krv:tbl1} приведены параметры, необходимые для настройки КРВ.

\footnotesize
\begin{longtable}{|>{\centering\arraybackslash}m{5.3cm}|>{\centering\arraybackslash}m{3.3cm}|>{\centering\arraybackslash}m{4.2cm}|>{\centering\arraybackslash}m{1.8cm}|>{\centering\arraybackslash}m{1cm}|}
\caption{Параметры для настройки КРВ\hfill\vspace{-0.5\baselineskip}}\label{krv:tbl1}\\ 
\hline
\rowcolor{gray!20} Параметр (Параметр на ИЧМ) & Условное обозначение на схеме & Значение/ Диапазон & Единица измерения & Шаг \\ 
\hline
\endfirsthead
\caption*{\hspace{3pt}\emph{Продолжение таблицы \ref{krv:tbl1}\hfill\vspace{-0.5\baselineskip}}} \\ % сделано по ГОСТ 2.105 п.6.8.7
\hline
\rowcolor{gray!20} Параметр (Параметр на ИЧМ) & Условное обозначение на схеме & Значение/ Диапазон & Единица измерения & Шаг \\ 
%\hline
\endhead
%\hline
\endfoot
%\hline
\endlastfoot 
%==+t1*CLS|CLS
\centering Ввод функции в работу (Ввод\_функции)  & \centering SGF1 & \centering Не предусмотрено \\ Предусмотрено & \centering -- & \centering \arraybackslash -- \\
\hline
\centering Минимальный ток КРВ (Iмин\_КРВ)  & \centering -- & \centering 0,01...0,30 & \centering о.е. & \centering \arraybackslash 0,01 \\
\hline
\centering Номинальный ток выключателя (паспорт) (Iном\_В\_пасп)  & \centering -- & \centering 0...5000 & \centering А & \centering \arraybackslash 1 \\
\hline
\centering Номинальный ток отключения выключателя (паспорт) (Iном\_откл\_В\_пасп)  & \centering -- & \centering 0...40000 & \centering А & \centering \arraybackslash 1 \\
\hline
\centering Допустимое количество циклов при номинальном токе выключателя (паспорт) (КРВпасп\_Iном)  & \centering -- & \centering 0...200000 & \centering -- & \centering \arraybackslash 1 \\
\hline
\centering Допустимое количество циклов при номинальном токе отключения выключателя (паспорт) (КРВпасп\_Iном\_откл)  & \centering -- & \centering 0...500 & \centering -- & \centering \arraybackslash 1 \\
\hline
\centering Допустимый механический ресурс выключателя (паспорт) (МРВпасп)  & \centering -- & \centering 0...200000 & \centering -- & \centering \arraybackslash 1 \\
\hline
\centering Начальное значение коммутационного ресурса выключателя (Нач\_знач\_КРВ)  & \centering -- & \centering 0,0...100,0 & \centering \% & \centering \arraybackslash 0,1 \\
\hline
\centering Величина срабатывания коммутационного ресурса выключателя (КРВсраб)  & \centering -- & \centering 0,0...100,0 & \centering \% & \centering \arraybackslash 0,1 \\
\hline
\centering Начальное значение механического ресурса выключателя (Нач\_знач\_МРВ)  & \centering -- & \centering 0...200000 & \centering -- & \centering \arraybackslash 1 \\
\hline
\centering Контроль КРВ по превышению МРВ (Контроль\_КРВ\_от\_МРВ)  & \centering SGF2 & \centering Не предусмотрено \\ Предусмотрено & \centering -- & \centering \arraybackslash -- \\
\hline
\centering Полное время отключения выключателя (Tмакс\_В)  & \centering -- & \centering 10...3000 & \centering мс & \centering \arraybackslash 10 \\
\hline
%===t1
\end{longtable}
\normalsize


%%%%%%%%%%%%%%%%%%%%%%%%%%%%%%%%%%%%%%%%%%%%%%%%%%%%%%%%%% ПС %%%%%%%%%%%%%%%%%%%%%%%%%%%%%%%%%%%%%%%%%%%%%%%%%%%%%%%%%%%%%%%%%%
\par
В таблице \ref{ps:tbl1} приведены параметры, необходимые для настройки ПС. 

\footnotesize
\begin{longtable}{|>{\centering\arraybackslash}m{5.3cm}|>{\centering\arraybackslash}m{3.3cm}|>{\centering\arraybackslash}m{4.2cm}|>{\centering\arraybackslash}m{1.8cm}|>{\centering\arraybackslash}m{1cm}|}
\caption{Параметры для настройки предупредительной сигнализации\hfill\vspace{-0.5\baselineskip}}\label{ps:tbl1}\\ 
\hline
\rowcolor{gray!20}
Параметр (Параметр на ИЧМ) & Условное обозначение на схеме & Значение/ Диапазон & Единица измерения & Шаг \\ 
\hline
\endfirsthead
%\caption{\emph{Продолжение\hfill\vspace{-0.5\baselineskip}}} \\ % Отступ обозначения таблицы от таблицы и выравнивание надписи слева
\caption*{\hspace{3pt}\emph{Продолжение таблицы \ref{ps:tbl1}\hfill\vspace{-0.5\baselineskip}}} \\ % сделано по ГОСТ 2.105 п.6.8.7
\hline
\rowcolor{gray!20}
Параметр (Параметр на ИЧМ) & Условное обозначение на схеме & Значение/ Диапазон & Единица измерения & Шаг \\ 
%\hline
\endhead
%\hline
\endfoot
%\hline
\endlastfoot
%==+t1*CALH1|LVALH
\centering Контроль сигнализации ТЗ (Контр\_ТЗ\_сигн)  & \centering SGF1 & \centering Не предусмотрено \\ Предусмотрено & \centering -- & \centering \arraybackslash -- \\
\hline
\centering Контроль сигнализации ГЗ (Контр\_ГЗ\_сигн)  & \centering SGF2 & \centering Не предусмотрено \\ Предусмотрено & \centering -- & \centering \arraybackslash -- \\
\hline
\centering Контроль низкой изоляции ТЗ (Контр\_низк\_изол\_ТЗ)  & \centering SGF3 & \centering Не предусмотрено \\ Предусмотрено & \centering -- & \centering \arraybackslash -- \\
\hline
\centering Контроль низкой изоляции ГЗ (Контр\_низк\_изол\_ГЗ)  & \centering SGF4 & \centering Не предусмотрено \\ Предусмотрено & \centering -- & \centering \arraybackslash -- \\
\hline
\centering Контроль блокировки ТЗ (Контр\_ТЗ\_блок)  & \centering SGF5 & \centering Не предусмотрено \\ Предусмотрено & \centering -- & \centering \arraybackslash -- \\
\hline
\centering Контроль блокировки ГЗ (Контр\_ГЗ\_блок)  & \centering SGF6 & \centering Не предусмотрено \\ Предусмотрено & \centering -- & \centering \arraybackslash -- \\
\hline
\centering Контроль выходных цепей (Контр\_цепей)  & \centering SGF7 & \centering Не предусмотрено \\ Предусмотрено & \centering -- & \centering \arraybackslash -- \\
\hline
\centering Контроль выведенного положения БИ (Контр\_вывода\_БИ)  & \centering SGF8 & \centering Не предусмотрено \\ Предусмотрено & \centering -- & \centering \arraybackslash -- \\
\hline
\centering Контроль сигнализации опер.тока (Контр\_ОТ\_сигн)  & \centering SGF9 & \centering Не предусмотрено \\ Предусмотрено & \centering -- & \centering \arraybackslash -- \\
\hline
\centering Контроль неисправности КА (Контр\_Неиспр\_КА)  & \centering SGF10 & \centering Не предусмотрено \\ Предусмотрено & \centering -- & \centering \arraybackslash -- \\
\hline
\centering Контроль превышения времени переключения КА (Контр\_прев\_вр\_пер)  & \centering SGF11 & \centering Не предусмотрено \\ Предусмотрено & \centering -- & \centering \arraybackslash -- \\
\hline
\centering Контроль общего внешнего сигнала (Контр\_общ\_вн\_сигн)  & \centering SGF12 & \centering Не предусмотрено \\ Предусмотрено & \centering -- & \centering \arraybackslash -- \\
\hline
\centering Контроль положения переключателя аварийной деблокировки (Контр\_пол\_ав\_дебл)  & \centering SGF13 & \centering Не предусмотрено \\ Предусмотрено & \centering -- & \centering \arraybackslash -- \\
\hline
%===t1
\end{longtable}
\normalsize

%%%%%%%%%%%%%%%%%%%%%%%%%%%%%%%%%%%%%%%%%%%%%%%%%%%%%%%%%% СС %%%%%%%%%%%%%%%%%%%%%%%%%%%%%%%%%%%%%%%%%%%%%%%%%%%%%%%%%%%%%%%%%%
\par
В таблице \ref{ss:tbl1} приведены параметры, необходимые для настройки СС. 

\footnotesize
\begin{longtable}{|>{\centering\arraybackslash}m{5.3cm}|>{\centering\arraybackslash}m{3.3cm}|>{\centering\arraybackslash}m{4.2cm}|>{\centering\arraybackslash}m{1.8cm}|>{\centering\arraybackslash}m{1cm}|}
\caption{Параметры для настройки функции сборки сигналов\hfill\vspace{-0.5\baselineskip}}\label{ss:tbl1}\\ 
\hline
\rowcolor{gray!20}
Параметр (Параметр на ИЧМ) & Условное обозначение на схеме & Значение/ Диапазон & Единица измерения & Шаг \\ 
\hline
\endfirsthead
\caption*{\hspace{3pt}\emph{Продолжение таблицы \ref{ss:tbl1}\hfill\vspace{-0.5\baselineskip}}} \\ % сделано по ГОСТ 2.105 п.6.8.7
\hline
\rowcolor{gray!20}
Параметр (Параметр на ИЧМ) & Условное обозначение на схеме & Значение/ Диапазон & Единица измерения & Шаг \\ 
%\hline
\endhead
%\hline
\endfoot
%\hline
\endlastfoot
%==+t1*GAPC1|OL_SignAssembly
\multicolumn{5}{|c|}{ CC } \\ \hline 
\centering Контроль опер. тока цепей ГЗ, ТЗ (Контр\_ОТ\_ГЗ\_ТЗ)  & \centering SGF1 & \centering Не предусмотрено \\ Предусмотрено & \centering -- & \centering \arraybackslash -- \\
\hline
\centering Контроль опер. тока цепей сигнализации выключателя (Контр\_ОТ\_сигн\_В)  & \centering SGF2 & \centering Не предусмотрено \\ Предусмотрено & \centering -- & \centering \arraybackslash -- \\
\hline
\centering Контроль опер. тока цепей ОБР (Контр\_ОТ\_ОБР)  & \centering SGF3 & \centering Не предусмотрено \\ Предусмотрено & \centering -- & \centering \arraybackslash -- \\
\hline
\centering Контроль положения SA1 (Контр\_полож\_SA1)  & \centering SGF4 & \centering Не предусмотрено \\ Предусмотрено & \centering -- & \centering \arraybackslash -- \\
\hline
\centering Контроль положения SA2 (Контр\_полож\_SA2)  & \centering SGF5 & \centering Не предусмотрено \\ Предусмотрено & \centering -- & \centering \arraybackslash -- \\
\hline
\centering Контроль положения SA3 (Контр\_полож\_SA3)  & \centering SGF6 & \centering Не предусмотрено \\ Предусмотрено & \centering -- & \centering \arraybackslash -- \\
\hline
\centering Контроль положения SA4 (Контр\_полож\_SA4)  & \centering SGF7 & \centering Не предусмотрено \\ Предусмотрено & \centering -- & \centering \arraybackslash -- \\
\hline
\centering Контроль положения SG1 (Контр\_полож\_SG1)  & \centering SGF8 & \centering Не предусмотрено \\ Предусмотрено & \centering -- & \centering \arraybackslash -- \\
\hline
\centering Контроль положения SG2 (Контр\_полож\_SG2)  & \centering SGF9 & \centering Не предусмотрено \\ Предусмотрено & \centering -- & \centering \arraybackslash -- \\
\hline
\centering Контроль положения SG3 (Контр\_полож\_SG3)  & \centering SGF10 & \centering Не предусмотрено \\ Предусмотрено & \centering -- & \centering \arraybackslash -- \\
\hline
\centering Контроль положения SG4 (Контр\_полож\_SG4)  & \centering SGF11 & \centering Не предусмотрено \\ Предусмотрено & \centering -- & \centering \arraybackslash -- \\
\hline
%===t1
\end{longtable}
\normalsize